% !TeX spellcheck = en_US
\documentclass[]{jnsao}
\usepackage{lineno}\linenumbers

\let\Bbbk\relax
\usepackage{amsmath,amsthm,amssymb}
\usepackage[a4paper, total={6in, 9in}]{geometry}
\usepackage{paralist}
\usepackage[T1]{fontenc}
%\usepackage{mathpazo} % Palatino font for text and math
\usepackage{microtype} % Improved typography
\usepackage[title,titletoc,page]{appendix}
\usepackage{hyperref}
%\DeclareMathOperator{\Lip}{Lip}
\usepackage{enumitem}
\usepackage{macros_perso}
%\newtheorem{proposition}{Proposition}
%\newtheorem{corollary}{Corollary}
%opening
\title{A new comparison principle for discrete Volterra equations with an application to convex sweeping processes}
\shorttitle{Sweeping processes with inifinite delays}
\author{Thierno Mamadou Baldé \thanks{Laboratoire de Mathématiques de Bretagne Atlantique,
Univ Brest, CNRS UMR 6205,
6, avenue Victor Le Gorgeu,
29200 Brest, France}
\and Vuk Milisic \footnotemark[1]
\and Steffen Plunder \thanks{Institute for Advanced Study of Human Biology (WPI-ASHBi), Kyoto University, Kyoto, 606-8303, Japan}
}
\shortauthor{T. M. Baldé, V. Milisic, S. Plunder}
\newcommand{\kernel}{\varrho}
\newcommand{\tkernel}{\tilde{\varrho}}
\newcommand{\tp}{\tilde{p}}
%\newcommand{\nrm}[2]{\|#1\|_{#2}}
\renewcommand{\tr}{R}
\newcommand{\tP}{\tilde{P}}
\newcommand{\tf}{\tilde{f}}
\DeclareMathOperator*{\argmin}{argmin}
\DeclareMathOperator*{\ret}{ret}

\begin{document}

\maketitle

\begin{abstract}
Comparison principles for Volterra equations play a role analogous to maximum principles in PDEs: 
they provide positivity and stability information on the solution and allow to control 
the output of bounded inputs. In the continuous setting, such results often rely on Laplace-transform 
or spectral methods (see e.g. Gripenberg et al. [1, Ch. 4. \& 7.]). However, these tools are not 
uniform in the discretization step  $h$ hence fail in discrete or semi-discrete approximations. 
The present note introduces a resolvent-free argument yielding uniform $L^\infty(0,T)$-bounds
for non-negative kernels.

Compactness is a key ingredient in order to show existence of sweeping processes.
While in the classical framework it is well established, adding an infinite 
distribution of delays complicates greatly the obtaintion of such a result.
In a first step we show a general energy decay estimate, which is then used to
establish compactness. The argument is carried out in the discrete setting
and that necessitates the introduction of the  new comparison principle.

In the classical sweeping process the previous position of the particle
lies on the boundary of the constraint set, staying $O(h)$
close to the next projection point ($h$ is the discretisation step). Our delay model projects the particle's
averaged (by a unit measure kernel) past positions  to the constraint set.
Numerical simulations show that the projected point can lie at $O(1)$
distance from the convex set's boundary.
\end{abstract}
%\section{Introduction}


\tableofcontents

\section{Introduction}

Volterra equations arise naturally in the modeling of systems with memory, where the present state depends
on the past history through a convolution kernel. In this work, we provide tools for analyzing a novel type of Volterra 
equations which is motivated by a problem from modelling of multi-celluar systems with adhesive memory and non-overlap 
conditions (see \cite{venel08} for the classical differential inclusion setting). 

Cells with adhesive memory have been modelled extensively using Volterra-type equations where the memory 
represents the forces of adhesion to past positions; see e.g.~\cite{OelzSch10,PreVi,Manhart2015,MiOel.1,MiOel.2,MiOel.3,MiOel.4} 
%\task{add also third party extensions}.
The system modelling the quasi-steady state of cells at each point in time reads 
\begin{align*}
  z(t) - \int_0^t k(t-s) z(s)\,ds = f(t), \quad t > 0,
\end{align*}
where $z(t)$ represents the position of the cell at time $t$, $k$ is a 
nonnegative kernel describing the adhesive memory, and $f(t)$ is an external force acting on the cell.

When multiple cells are considered, one might want to impose non-overlapping constraints between cells, effectively constraining the state to a feasible set $C(t)$ at each time $t$. 
This leads to a new Volterra-type sweeping process: find $z(t) \in C(t)$ for $t>0$, such that
$$
z(t) - \int_0^t k(t-s) z(s)\,ds - f(t) \in -N_{C(t)}(z(t)),\quad t > 0,
$$
where $C(t)$ can be for instance the 
interior convex approximation of the non overlapping constraint 
and simplified to a time-dependent convex set and $N_{C(t)}(z(t))$ is the normal cone to $C(t)$ at $z(t)$.

As the movement of the set $C(t)$ might be much faster than the natural dynamics of the unconstrained 
system, we experience here a challenging geometric setup. This becomes visible when the dynamics are 
recast as a projective system, such as 
\begin{align*}
  z(t) = P_{C(t)}\Big( \int_0^t k(t-s) z(s)\,ds + f(t) \Big), \quad t > 0.
\end{align*}
Here the projection distance depends crucially on the kernel and on the past trajectory $z(t)$.
This shows the fundamental difference with  classical sweeping processes. Indeed, for Moreau's model,
  the projection distance is infinitesimal due to the presence of a time-derivative term.
Dealing with non-infinitesimal projection distances requires new bounds on the solution $z(t)$ in order to 
establish stability and compactness. Often the stability result 
relie on some kind of Gronwall's lemma (for odes for instance) or more generally on comparison principles.
For Volterra equations these are not straightforward to obtain, this is due to the non-local nature of the convolution operator
and from the physical perspective, the memory effect implies that the solution at time $t$ depends on the entire past history,
and can exhibit oscillations or growth that are not present in local differential equations.
Comparison principles for such equations play a 
role analogous to maximum principles in parabolic PDEs: they ensure that bounded or positive data produce 
bounded or positive solutions, and they provide essential \emph{a~priori} bounds for qualitative and numerical analysis.
 
In the continuous setting, comparison results can often be obtained by characterizing 
the resolvent operator associated with the Volterra kernel through the resolvent equation
\[r(t) - \int_0^t \kernel	(t-s) r(s)\,ds = \kernel(t), \qquad t > 0,\]
then using the Laplace transform to analyze solutions' behavior provides a spectral
decomposition of the resolvent in a sum of exponentials and a remainder term that
is often integrable or decays in a sufficient manner \cite[Ch.~4 \& 7]{grip}. 
However, in many applications—especially when the kernel is not exponentially decaying or when it depends on 
the discretization parameter—such a spectral decomposition is not directly available. Moreover, 
even if Wiener’s lemma can be invoked to provide a similar asymptotic description of the discrete resolvent, 
the corresponding representation is not uniform with respect to the discretization step $h$ 
and therefore does not provide uniform \emph{a~priori} bounds for the discrete solutions.






%\medskip
%\noindent
%\textbf{Motivation and contribution.}
Classical convergence and stability analyses of discrete Volterra equations, 
such as those of McKee and Wall~\cite{JonesMcKee1982}, Hickernell~\cite{BakkeJackiewicz1988}, 
and Lubich~\cite{Lu.88.1,Lu.88.2}, rely on spectral or Laplace-transform techniques that 
require exponentially decaying or  kernels whose Laplace transform is well-behaved. 
These approaches provide powerful tools 
for these kernels but fail to yield uniform bounds when the memory kernel is not explicit 
or lacks analyticity. The present work develops a comparison principle that 
circumvents the use of any spectral decomposition and applies equally to continuous and 
discrete settings. The key novelty lies in the construction of an {\em initial layer} 
corrector compensating for the long-time tail of the kernel, which ensures positivity 
and boundedness of the discrete solution uniformly with respect to the discretization step. 


The comparison principle found in~\cite[Theorem~2.7]{MiOel.1} does not hold in the case of a 
non-exponential kernel that is constant in time. Here we construct a method that avoids the 
use of the spectral decomposition of the resolvent, which could be unavailable or $h$-dependent
without any uniform control in $h$ (if $h$ is the discretisation step-size for instance). 
Our goal is to develop a comparison argument that remains stable when 
the equation is discretized in time and that yields uniform bounds for the corresponding 
discrete solutions. This approach is motivated by applications to viscoelasticity and 
adhesion models, where the memory kernel may decay only algebraically and where time 
discretizations must preserve the qualitative properties of the continuous dynamics.




Moreau's sweeping process (classical formulation) 
is the baseline model: a first-order differential inclusion 
  $-\dot X(t)\in N_{C(t)}(X(t))$ introduced by J.-J. Moreau \cite{Moreau1999}.
  Delay and history-dependent perturbations have been studied since: 
  existence results for sweeping processes with delays and history operators were 
  established in works by Castaing \& Monteiro Marques \cite{Castaing1997} and 
  Edmond \cite{Edmond2006}.
 Most results for delayed sweeping processes keep a time-derivative
 term (i.e. explicit $\dot X$ or its fractional analogue) in the formulation; 
 the analysis relies on this differential feature to obtain compactness and pass-to-the-limit.
Fractional-time derivative sweeping processes (Caputo-type) have been developed recently; 
  see e.g. Zeng \cite{Zeng2023} and Bouach \cite{Bouach2025}, 
  by a fractional derivative and yield existence/uniqueness 
   and—importantly—compactness properties coming from fractional regularizing effects.
By contrast, replacing the time-derivative by a general Volterra (integro-differential / history) 
operator leads to different analytic difficulties: compactness is generally lost unless very 
strong hypotheses are imposed. 
Recent papers prove well-posedness for Volterra-type integro-differential 
sweeping processes using tailored a priori estimates; 
see Vilches \cite{Vilches2024}, Godoy et al. \cite{Godoy2024}, and Haddad \cite{Haddad2025}.
State-dependent (moving set depends on $x$) and prox-regular (non-convex) extensions have been 
developed for both delay and fractional setups; proofs typically combine prox-regular geometry 
with fixed-point or reparametrization techniques.
Numerical and applied analyses (contact, frictionless impact, plasticity) 
motivate the integro-differential and fractional variants.

Our work contributes to this recent line of research by establishing a new comparison principle
for discrete integral Volterra equations with non-exponential non-negative kernels. 
To be more precise, our problem does not involve a time-derivative term,
rather it is replaced by a memory operator involving an infinite distribution of delays.
This operator does has been introduced in \cite{OelzSch10} to account
for a microscopic friction mechanism mediated by transient elastic linkages to past positions.
Biological considerations motivate this model, as cells adhere to their past positions
through transient bonds that form and break over time.
In order to show existence of sweeping processes relying on  such principles, 
we use minimizing movements à la De Giorgi. The key step is to obtain
uniform $H^1$-bounds on the discrete solutions, which then provides
compactness by Ascoli-Arzelà theorem. The new comparison principle
allows to obtain such uniform bounds in the discrete setting.
This is crucial for passing to the limit and establishing existence of solutions.
To our knowledge, there is no such a result in the literature yet.
Even for the case of the simplest sweeping process where the moving convex set 
is a sphere.






The paper is organized as follows. 
In Section \ref{sec:assumptions}, we gather all the assumptions and state the main results.
Section~\ref{sec:continuous} recalls the 
continuous framework and basic assumptions ensuring the existence of a unique solution. 
In Section~\ref{sec:initiallayer}, we study a canonical auxiliary problem that reveals 
the structure of the \emph{initial layer} associated with the convolution kernel. 
Building on this analysis, Section~\ref{sec:comparison} establishes a new comparison 
principle valid for non-exponential kernels, without resorting to a spectral decomposition. 
Section~\ref{sec:discrete} introduces a discrete counterpart of the continuous model, 
derives a discrete conservation relation, and proves uniform $L^\infty$ bounds independent 
of the discretization step $h$. Section~\ref{sec:application} applies these results to a 
delayed sweeping process, establishing energy estimates and compactness properties needed 
for existence results. Then by standard arguments one passes to the limit in the discrete
approximations and recovers an existence result for the continuous delayed sweeping process.
Finally, the Appendices \ref{sec:moments},  \ref{sec:discrete_convolution} and \ref{sec:decay_lipschitz} contain auxiliary technical results.


\section{Assumptions and main results}\label{sec:assumptions}

In this section, we gather the assumptions used 
throughout the paper and state the main results.
Assumptions concerning the Volterra equations that
we consider are stated hereafter.
\begin{hypo}[Data assumptions]\label{hypo.data}
Throughout this paper, we assume the following conditions on the data:
	\begin{compactenum}[i)]
		\item the source term  $f \in L^\infty (\rr)$,  
		\item the kernel $\kernel$ is non-negative unit measure and satisfies $\kernel \in L^1(\rr,(1+a)^2)\cap \Lip(\rr)$,
		\item the kernel satisfies the monotonicity condition: $\da \kernel (a) \leq 0$,
		\item there exists a constant $\zeta > 0$ such that $- \da \kernel(a) / \kernel(a) \leq \zeta$ for all $a>0$,
		\item the past condition reads: $\zp \in \Lip (\RR^-)$.
	\end{compactenum}
\end{hypo}

In the rest of the paper the renormalization of the kernel is made for 
sake of conciseness and all the results hold for generic non negative
integrable kernels.
\renewcommand{\tkernel}{\kernel}


\noindent For the sweeping process, we require some regularity on the 
moving convex sets:
\begin{hypo}[Convex sets assumptions]\label{hypo.convex.main}
For the application to delayed sweeping processes, we consider a family of 
closed convex sets $(C^n)_{n\in\N}$ in $\mathbb{R}^d$, satisfying:
\begin{compactenum}[i)]
	\item  there exists $R>0$ such that for all $n\in\N$, $C^n \subset B(0,R)$,
	\item compactness of the sequence: for all $h<h_0$, 
	setting $N=\lfloor T/h \rfloor$, the sequence of convex sets $(C^n)_{0\leq n \leq N}$ satisfies the $H^1$-like estimate:
	\begin{equation}\label{eq.bv.c.main}
\sum_{j=1}^N d_H(C^{n},C^{n-1})^2/h \leq C < \infty,
\end{equation}
where the constant $C$ is uniform with respect to $h$,
%\item Lipschitz continuity: $d_H(C^{n},C^{n-1}) \leq C h$ for all $n\geq 1$.
\end{compactenum}
\end{hypo}

%\subsection{Main results}
\subsection{Comparison principles for Volterra equations: continuous and discrete settings}
In a first step, and for sake of clarity, we show the 
theoretical construction leading to the comparison principle 
in the continuous setting.
\begin{theorem}[Comparison principle for continuous Volterra equations]\label{thm.main.comparison}
	Assume that $z$ solves:
	\begin{equation}\label{eq.z.main}
		\left\{\begin{aligned}
			&	 \int_\rr (z(t)-z(t- a)) \kernel(a) da = f(t), & t>0 \\
			&	z(t) = \zp(t), & t<0,
		\end{aligned}\right.
	\end{equation}
	under Hypothesis \ref{hypo.data}. Then $z \in L^\infty (0,T)$ for any finite $T$, and the bound is independent of the spectral decomposition of the resolvent.
\end{theorem}
This result is made possible introducing an initial layer corrector.
For sake of clarity,  technical aspects are further detailed in section \ref{sec:initiallayer}. %as follows:
% \begin{thm[Boundedness of initial layer]\label{thm.w.main}
% 	Let $w$ be the solution of the initial layer problem:
% 	\begin{equation}\label{eq.w.main}
% 		 w(t)-\int_0^t w(t-a) \tkernel(a) da = \int_t^\infty \left(a-t \right) \tkernel(a) da, \;  t>0,
% 	\end{equation}
% 	with $w(t) = -t$ for $t< 0$, where $\tkernel (a):= \kernel(a)/ \int_\rr \kernel(\tia) d\tia$.
% 	Under Hypothesis \ref{hypo.data}, we have $w \in L^\infty (\rr)$ and $w \geq 0$ for a.e. $t\geq 0$.
% \end{theorem}
\noindent Mimicking the continuous setting, we now turn to the discrete framework.
We discretize the time interval $[0,T]$ with a step $h>0$ and
 we denote $t_n:= n h$ for all $n \in \{0, \ldots, N\}$ where $N:= T/h$. 
 We assume for sake of simplicity that $T/h \in \N$.
 Then we use piecewise constant approximations. 
 Whithin this framework, we state the following corresponding result.
\begin{theorem}[Uniform discrete comparison principle]\label{thm.comp.disc.main}
	Consider the discrete problem:
	\begin{equation}\label{eq.disc.z.main}
	\left\{\begin{aligned}
		&	Z^n - h \sum_{j\geq 0} Z^{n-j} \tr_j = f^n , & n \in \{0, \ldots, N\} \\
		&	Z^n = Z_p^n , & n<0,
	\end{aligned}\right.
	\end{equation}
	where $R_j:= \frac{1}{h} \int_{jh}^{(j+1)h} \kernel(a) da$.
	Assume that %$-(R_j-R_{j-1})/(hR_j) \leq \zeta$ for all $h>0$ and all $j\geq 1$, 
	Hypotheses \ref{hypo.data} holds. Then the solution $(Z^n)_{n\in \{0,\dots,N\}}$ of \eqref{eq.disc.z.main} is bounded uniformly with respect to $h$.
\end{theorem}
\noindent Again the key ingredient is the discrete initial layer corrector and its 
control (cf section \ref{subsec:discrete}).
% \begin{proposition[Discrete initial layer]\label{prop.W.bounded.main}
% 	The discrete initial layer $(W^n)_{n\in\N}$ solving:
% 	\begin{equation}\label{eq.disc.w.main}
% 	\left\{\begin{aligned}
% 		&	 W^n - h\sum_{j\geq 1} W^{n-j} \tr_j =0 & n \geq 0 \\
% 		& W^n = - nh & n< 0,
% 	\end{aligned}
% 	\right.
% 	\end{equation} 
% 	is a bounded non-negative sequence, with bound uniform in $h$.
% \end{proposition}

\subsection{Energy estimates for delayed sweeping processes}

We formulate our sweeping process as the projection of the averaged 
past position onto the convex set:
$$
X^n:= P_{C^n} (\overline{X}^n), \quad 
\overline{X}^n :=  \frac{h \sum_{j\geq 1} X^{n-j} \tr_j}{h \sum_{j\geq 1} \tr_j}, \quad n \in \{0, \ldots, N\}.	
$$
where the kernel's weights $(\tr_j)_{j\in \N}$ are defined as above.
Then we show that this sweeping process can be equivalently reformulated as the following
delayed gradient flow \cite{Mi.5} and that it satisfies the energy dissipation estimates:
\begin{proposition}[Energy dissipation]\label{prop.nrj.main}
\begin{equation}\label{eq.sweeping.main}
	 X^n := \argmin_{W \in C^n} \cE_n(W), \quad \cE_n(W) := \frac{h}{2}\sum_{j\geq1} (W-X^{n-j})^2 \tr_j.
	\end{equation}
	Assume that $(\tr_j)_{j\in \N}$ is decreasing and Hypothesis \ref{hypo.convex.main} holds. 
	Then:	
	\[
	\cE_{N}(X^{N}) +  \sum_{n=0}^{N-1} D_n \le e^{2 T} \cE_{0}(X^0) + e^{2 T} \sum_{n=0}^{N-1} e(C^n,C^{n+1})^2/h,
	\]
	where $D_n= \frac{h}{2} \sum_{j\ge1} (\tr_{j}-\tr_{j+1}) | X^{n} - X^{n-j}|^2 \geq 0$ is the dissipation term.
\end{proposition}

% \begin{corollary}[Dissipation estimate]\label{coro.dissip.main}
% 	Under the hypotheses of Proposition \ref{prop.nrj.main},
% 	$$
% 	\sum_{p= 0}^N h 
% 	\left( \sum_{j\ge1} (\tr_{j-1} - \tr_{j}) \left| X^{p} - X^{p-j}\right| \right)^2 \leq  C < \infty, 
% 	$$
% 	uniformly with respect to $h$.
% \end{corollary}

% \begin{corollary}[Lagrangian bound]\label{coro.lag.main}
% 	Under the hypotheses of Proposition \ref{prop.nrj.main}, the distance between the averaged past and the current position is uniformly bounded:
% 	$$
% 	\left|\overline{X}^n - X^n\right| \leq C, \quad \text{where } \overline{X}^n := h \sum_{j\geq 1} X^{n-j} \tr_j.
% 	$$
% \end{corollary}

\subsection{Convergence to continuous delayed sweeping process}
Define the piecewise linear interpolant $X_h$ of the discrete solution $X^n$ on $(0,T)$.
In the same way, define the piecewise constant interpolant $\overline{X}_h$ of the discrete averages $\overline{X}^n$ on $(0,T)$.
Then, we have the following convergence result. 
\begin{theorem}[Limit of the discrete sweeping process]\label{thm:limit_sweeping.main}
	Under Hypotheses \ref{hypo.data} and \ref{hypo.convex.main}, assume the uniform estimate
	$ \| X_h \|_{H^1(0,T)}\le C$ for all $h>0$, for any finite $T>0$.
	
	Then, as $h\to0$, there exist $X\in H^1([0,T];\mathbb R^m)$ and a subsequence such that 
	$X_h \to X$ in $C^0 (0,T;\mathbb R^m)$, and the discrete averages converge 
	$\overline{X}_h \to \overline{X}$ in $L^p(0,T;\mathbb R^m)$ for $1\le p<\infty$, where
	\[
	\overline{X}(t) := \int_0^\infty \varrho(s)\,X(t-s)\,ds.
	\]
	Moreover, the limit function $X$ satisfies the Moreau-type integral inclusion:
	\begin{equation}\label{eq:limit_inclusion.main}
		\overline{X}(t)-X(t)\in N_{C(t)}\big(X(t)\big), \quad\text{for a.e. }t\in(0,T),
	\end{equation}
	equivalently, $X(t)=P_{C(t)}\big(\overline{X}(t)\big)$ for a.e. $t\in(0,T)$.
\end{theorem}

\begin{remark}
	The general convergence result stated in Theorem \ref{thm:limit_sweeping.main} 
	requires an $H^1(0,T)$ bound on the discrete solution, which is difficult to 
	establish for general convex sets. The following result provides such a bound 
	for the special case of moving circular sets. In order to obtain these 
	uniform $H^1$ bounds, we crucially rely on the new comparison principle
	presented above.
\end{remark}

\begin{theorem}[Convergence for moving circular sets]\label{thm:circular.main}
	Consider the case where $C(t) := \overline{B}(c(t), R)$ with $c \in H^1([0,T];\mathbb{R}^m)$ and $R \in \mathbb{R}_+$.
	Discretize the sets as $C^n := C(nh)$ for $n=0,\dots,N$ with $N = \lfloor T/h \rfloor$.
	Assume that $(\tr_j)_{j\geq 0}$ is non-increasing with mass one, and that $X^{-j} = X^{-1} \in C_0$ for all $j\geq 2$.
	
	Then the discrete sweeping sequence $(X^n)_{n\ge0}$ defined by \eqref{eq.sweeping} satisfies the uniform bounds:
	\[
	\sup_{n\in \{0,\dots,N_T\}}|X^n|\le C,\qquad
	\left\| \dt X_h \right\|^2_2 := \frac{1}{h}\sum_{n=0}^{N-1} |X^{n+1}-X^n|^2  \le C,
	\]
	where $X_{h}$ is the piecewise-linear interpolant of $X^n$ on $(0,T)$, and $C>0$ is independent of $h$.
	
	Consequently, by Theorem \ref{thm:limit_sweeping.main}, the discrete solutions converge to a continuous delayed sweeping process satisfying $X(t)=P_{C(t)}\big(\overline{X}(t)\big)$ for a.e. $t\in(0,T)$.
\end{theorem}


\section{The continuous setting}\label{sec:continuous} 

We start with the problem : find $z \in L^1_\loc(\rr)$ solving  
\begin{equation}\label{eq.z}
	\left\{\begin{aligned}
& 	\int_\rr (z(t)- z(t-a)) \kernel(a) da = f(t), & t>0 \\
& 	z(t) = \zp(t), & t<0,
	\end{aligned}\right.
\end{equation}
where the data satisfy Assumptions \ref{hypo.data}. 
%\begin{hypo}\label{hypo.data} the data satisfy the following hypotheses :
%	\begin{compactenum}[i)]
%		\item the source term  $f \in L^\infty (\rr)$,  
%		\item the kernel is non negative and satisfies $\kernel \in L^1(\rr,(1+a)^2)$ ,
%		\item the past condition reads : $\zp \in \Lip (\RR^-)$
%	\end{compactenum}
%\end{hypo}
By standard arguments :
\begin{theorem}
	Under assumptions \ref{hypo.data}, there exists a unique $z \in L^1_\loc(\rr)$ solving \eqref{eq.z}.
\end{theorem}
\noindent The proof is classical and relies on the construction of a positive resolvent and  can be found for instance in \cite[Theorem 3.5, p. 44]{grip}.

The solution can be represented by the 
formula $z(t) = f(t) + \int_0^t r(t-a) f(a) da$ 
where $r$ is the resolvent associated with $\kernel$.
Then techniques based on the Laplace transform lead to 
a spectral decomposition of the resolvent giving 
sharp $L^p$ estimates of the solution in a direct way.
Our aim is to provide $L^\infty$-bounds on $z$ that do not rely on 
the knowledge of $r$.
Indeed, we expose here ideas that should be efficient at the discrete level
where the spectral decomposition does not provide uniform asymptotic characterization
with respect to $h$. 
Our approach relies on the concept of a comparison principle.
For sake of clarity, we detail the basic ideas first in the continuous setting and then
transpose these to the discrete level.
For this purpose, we first study an initial layer corrector.

 




\subsection{An initial layer}\label{sec:initiallayer}
We solve for the time being 
a given problem whose importance will be explained later on. 
We denote by $\tkernel$ the renormalized version of the convolution kernel above:
$\tkernel (a) := \kernel(a)/ \int_\rr \kernel(\tia) d\tia$.
\begin{equation}\label{eq.w}
%%	\begin{aligned}
		 w(t)-\int_0^t w(t-a) \tkernel(a) da = \int_t^\infty \left(a-t \right) \tkernel(a) da, \;  t>0
%%	\end{aligned}
\end{equation}
This problem can also be  reformulated as 
\begin{equation}\label{eq.L.conserved}
	\left\{
	\begin{aligned}
		& \int_\rr (w(t)-w(t-a))\kernel(a) da = 0 & t\geq 0\\
		& w(t) = -t & t< 0\\
	\end{aligned}
	\right.
\end{equation}
Setting the elongation variable as $u(a,t) := w(t)-w(t-a)$, one can 
also reformulate the problem as a non-local PDE on this new variable 
\cite{MiOel.1}: find $u \in L^\infty((0,T)\times\rr,(1+a)^{-1}))$ solving 

\begin{equation}\label{eq.elongation.u}
	\left\{
\begin{aligned}
	& \dt u + \da u = - \int_\rr \da \tkernel(\tia) u(\tia,t)d \tia, & a>0,\; t>0\\
	& u(0,t)=0, & a=0,\; t>0\\
	& u(a,0)= w(0^+) - w(-a)=:u_I(a), & a>0,\; t=0
\end{aligned}
\right.
\end{equation}

Indeed, formally one has: $\dt u + \da u = \dot{w}$ by definition 
of $u$, then one writes:
\begin{align}
\ddt{} \int_\rr \tkernel(a) u(a,t) da &- \int_\rr \da \tkernel (a ) u(a,t) da = \left(\int_\rr \tkernel(a) da\right) \dot{w} = \dot{w}
\end{align}
but with this new definition \eqref{eq.L.conserved} becomes simply  
$\int_\rr \tkernel(a) u(a,t) da=0$ for a.e. $t>0$, this shows that 
$\dot{w} = - \int_\rr \da \tkernel(a) u(a,t) da$.  
Then using the first transport equation above show how, starting from $w$ solving \eqref{eq.w},
one recovers $u$ solving \eqref{eq.elongation.u}.
Conversely, starting from $u$ solving \eqref{eq.elongation.u}, one recovers $w$ solving \eqref{eq.w} by setting 
$w(t) := w(0^+) + \int_0^t \dot{w}(s) ds$ with $\dot{w}(s) = - \int_\rr \da \tkernel(a) u(a,s) da$. 
Although these computations (and the rest of formal computations of this section) 
are formal, they can be made rigorous using the 
method of characteristics (see for instance the detailed computations of \cite[Lemma 3.1]{MiOel.1}).
A simple computation shows that $w(0^+)= \mu_1/\mu_0 = \mu_1$ where 
$\mu_k := \int_\rr a^k \kernel(a) da$.
Using \cite[Theorem 6.1]{MiOel.2}, there exists a unique solution 
$u \in L^\infty((0,T)\times\rr,(1+a)^{-1}))$ of \eqref{eq.elongation.u}.

We are interested in a bound on $w$ so it seems natural to integrate 
\eqref{eq.elongation.u} in time. This leads to define 
$p(a,t) := \int_0^t u(a,t) da$ which solves:
\begin{equation}\label{eq.p}
	\left\{
	\begin{aligned}
		& \dt p + \da p= - \int_\rr \da \tkernel(\tia) p(\tia,t)d \tia + u_I(a),& a>0,\; t>0\\
		& p(0,t)=0 & a=0,\; t>0\\
		& p(a,0)= 0 & a>0,\; t=0
	\end{aligned}
	\right.
\end{equation}
we then make a change of variables to get rid of the $a$-dependent extra source term 
that appears on the right hand side. 
Setting $\tp(a,t) := p(a,t) - \int_0^a u_I(s)ds$, it solves:
 \begin{equation}\label{eq.tp}
 	\left\{
 	\begin{aligned}
 		& \dt \tp + \da \tp= - \int_\rr \da \tkernel(\tia) \tp(\tia,t)d \tia  - \int_\rr \da \tkernel  \int_0^a u_I(s) ds, & a>0,\; t>0\\
 		& \tp(0,t)=0 & a=0,\; t>0\\
 		& \tp(a,0)= - \int_0^a  u_I(s)ds & a>0,\; t=0
 	\end{aligned}
 	\right.
 \end{equation}
 Using then Proposition \ref{prop:decay_weighted}, the integration by parts formula is well defined and we have:
 $$
 \int_\rr \da \tkernel \int_0^a u_I(s)ds da = - \int_\rr \tkernel(a) u_I(a) da = 0.
 $$ 
 so that \eqref{eq.tp} reduces simply to: 
$$
\left\{
\begin{aligned}
	& \dt \tp + \da \tp= - \int_\rr \da \tkernel(\tia) \tp(\tia,t)d \tia ,& a>0,\; t>0\\
	& \tp(0,t)=0 & a=0,\; t>0\\
	& \tp(a,0)= - \int_0^a  u_I(s)ds & a>0,\; t=0
\end{aligned}
\right.
$$

\begin{theorem}\label{thm.w}
	Assume that $\da \kernel (a) \leq 0$ and that there exists a constant $\zeta > 0$ s.t. $- \da \kernel / \kernel < \zeta$, then 
	$w \in L^\infty (\rr)$ and $w \geq 0$ for a.e. $t\geq 0$.
\end{theorem}

\begin{proof}[Proof of Theorem \ref{thm.w}]
 	We apply the key idea of \cite[Lemma 5.1]{MiOel.2} and write:
 	\begin{align}
 	\ddt{} \int_\rr \tkernel (a) |\tp(a,t)| da &- \int_\rr \da \tkernel (a) |\tp(a,t) | da \leq - \int_\rr \da \tkernel (a) |\tp(a,t) | da
 	\end{align}
 	providing that 
	\begin{align}
	\int_\rr \tkernel(a) | \tp(a,t) | da &\leq \int_\rr \tkernel(a) | \tp(a,0) | da 
	= \int_\rr \tkernel(a) \left| \int_0^a u_I\right| da \lesssim \| \tkernel (1+a)^2\|_{1}.
	\end{align}
 	Then writing that 
 	$$
 	\begin{aligned}
 		w(t) = & w(0) + \int_0^t \dot{w}(s)ds = w(0) - \int_0^t \int_\rr \da \tkernel(a) u(a,s) da ds  
		= w(0) -  \int_\rr \da \tkernel(a) p(a,t) da, 
 	\end{aligned}
 	$$
 	this provides that:
 	\begin{align}
 			\left|w(t)\right| &\leq \frac{\mu_1}{\mu_0} + \int_\rr \frac{-\da \tkernel}{\tkernel} \tkernel(a) | p(a,t) | da\leq \frac{\mu_1}{\mu_0} + \zeta \int_\rr   \tkernel(a) | p(a,t) | da \\
 			&\leq C + \zeta \int_\rr | \tp(a,t) | \tkernel(a) da +  \zeta \int_\rr \left| \int_0^a u_I(s)ds \right| \tkernel(a) da \\
			&\leq C \left(1 + \| \tkernel \|_{L^1(\rr,(1+a)^2)}\right).
 	\end{align}
The initial layer corrector $w$ is thus bounded in $L^\infty(\rr)$. Moreover it is non-negative since the source term in \eqref{eq.w} is non-negative and the kernel is also non-negative.
Indeed, the resolvent associated with \eqref{eq.w} is a positive: 
as we assumed that $0\leq -\da \kernel(a) / \kernel(a) \leq \zeta$, 
 	% $$
 	% \int_t^\infty \kernel(a) da \geq \kernel(0) \int_t^\infty \exp(-\zeta a) da, \quad \forall t>0,
 	% $$
 	which gives immediately that 
 	$$
 	\int_0^t \tkernel (a) da = 1 - \int_t^\infty \tkernel(a) da \leq 1 - \frac{\kernel(0)\exp(-\zeta t)}{\zeta \mu_0} < 1.
 	$$
 	We are in the position to  apply \cite[Proposition 9.8.1]{grip}, which ensures the positivity of the resolvent. Next, one has the 
	representation formula for the solution $w$ of \eqref{eq.w}: $w(t)= \int_t^\infty(a-t)\tkernel(a) da + \int_0^t r(t-a) \int_a^\infty (s-a)\tkernel(s) ds da\geq 0$.
	\end{proof}
 
 \subsection{A new comparison principle}\label{sec:comparison} 

\begin{theorem}\label{thm.comparison}
	Assume that $z \in L^1_\loc(\rr)$ solves:
	\begin{equation}\label{eq.z.comparison}
		\left\{\begin{aligned}
			&	 \int_\rr (z(t)-z(t- a)) \kernel(a) da = f(t), & t>0 \\
			&	z(t) = \zp(t), & t<0,
		\end{aligned}\right.
	\end{equation}
	under Assumptions \ref{hypo.data}. 
	%Assume moreover that there exists $\zeta >0$ s.t. $-\da \kernel(a) / \kernel(a) \leq \zeta$ and that $\kernel$ is a 
	%non-increasing integrable function in the weighted $L^1(\rr,(1+a)^2)$ space 
	then $z \in L^\infty (0,T)$ for any finite $T$.   
\end{theorem}
 
 \begin{proof}[Proof of Theorem \ref{thm.main.comparison}]
 	Our aim is now to construct a super-solution of our problem, for this purpose we start as in the proof of \cite[Theorem 1.1]{MiOel.1} and we define :
 	\begin{equation}\label{eq.def.U}
 	U_f(t) :=  C  + \frac{\mu_0}{\mu_1} \begin{cases}
 		\int_0^t \nrm{f}{L^\infty(s,T)} ds & t > 0 \\
 		t \nrm{f}{L^\infty(0,T)}& \text{otherwise}
 	\end{cases}
 	\end{equation}
 	and we complete this profile by adding an initial layer, 
	defining $S_f(t) := U(t) +\frac{ \|f\|_\infty }{\mu_1} w(t)$, where we denote 
	$\|f\|_\infty := \|f\|_{L^\infty(0,T)}$. 
	We compute 
 	\begin{align}
 		&S_f(t) - \int_0^t S_f(t-a) \tkernel(a) da \\
 		&= \int_0^\infty \left(U_f(t) - U_f(t-a)\right) \tkernel (a) da  
		+ \int_t^\infty U_f(t-a) \tkernel(a) da 
		+\frac{ \|f\|_\infty  }{\mu_1}  \int_t^\infty (a-t) \tkernel(a) da,  
 	\end{align}
 	by the same arguments as in the proof of \cite[Theorem 1.1]{MiOel.1}, one writes:
	$$
	\begin{aligned}
	& \left(\int_0^t + \int_t^\infty \right)\left( U(t)-U(t-a)\right) \tkernel(a)da \\
	& \geq 
	\frac{\mu_0}{\mu_1} \left\{\int_0^t \|f\|_{L^\infty(t,T)} a \tkernel(a) da 
	+ \int_t^\infty t \|f\|_{L^\infty(t,T)}  + (a-t) \|f\|_{L^\infty(0,T)} \tkernel(a) da \right\}\geq |f(t)|. 
	\end{aligned}
	$$
	it remains to take care of the non-positive tail $\int_t^\infty U_f(t-a) \tkernel(a) da$ added in order to obtain the latter inequality.	
	one concludes that
 	\begin{align}
 	&S_f (t) - \int_0^t S_f(t-a) \tkernel(a) da \\
	&\geq | f(t) | + \int_t^\infty \left(C+ \frac{\nrm{f}{L^\infty(0,T)} }{\mu_1} (t-a) \right)\tkernel(a) da  +  \frac{ \|f\|_\infty  }{\mu_1}  \int_t^\infty (a-t) \tkernel(a) da  \geq | f(t) | 
 	\end{align}
 	the major improvement is here the compensation that our initial layer corrector $w$ provides: the linear part of $U$'s tail is negative and exactly 
 	matching the positive contribution from $w$. 
	As in the proof of Theorem \ref{thm.w},  the resolvent being positive one can apply \cite[Lemma 9.8.2 p. 257]{grip} in order to conclude. 
	Namely 
	$$
 	\begin{aligned}
 	|z(t)| - \int_0^t & |z(t-a)| \tkernel(a) da \\
	&  \leq \frac{|f(t)| }{\mu_0}
	+ \int_t^\infty \left|  \zp(t-a) \right|\tkernel(a) da =: \tilde{f} \leq S_{\tilde f} (t)- \int_0^t S_{\tilde f}(t-a) \tkernel(a) da 
 	\end{aligned}
 	$$
 	which implies that $|z(t)| \leq S_{\tilde{f}}(t)$, as soon as $|z(0)| \leq S(0)$, the latter condition is fulfilled by tuning the value of $C$ in the definition of $U$ in \eqref{eq.def.U}.
 \end{proof}
 
 \begin{remark}
 	The condition $-\da \kernel (a) / \kernel(a) \leq \zeta$  allows to consider polynomial decrease of the kernel at infinity. The method of recovering the tails introduced in \cite{MiOel.1}
 	in Lemma 2.4 fails for the latter kernels. If one considers only $U$ as a super-solution, then applying the integral operator implies:
 	$$
 	\begin{aligned}
 		U(t) - &\int_0^t U(t-a) \tkernel(a) da \geq |f(t)| + C \int_t^\infty \tkernel(a) da + \frac{\|f\|_\infty}{\mu_1} \int_t^\infty (t-a) \tkernel(a) da \\
 		& =  |f(t)| + \int_t^\infty \tkernel(a) da \left(C + \frac{\|f\|_\infty}{\mu_1} \frac{\int_t^\infty (t-a) \tkernel(a) da}{\int_t^\infty  \tkernel(a) da} \right)  \\
 			& = :
 			|f(t)| + \int_t^\infty \tkernel(a) da \left(C - \frac{\|f\|_\infty}{\mu_1} A [\tkernel](t) \right) 
 	\end{aligned}
 	$$
 	In \cite[Lemma 2.4]{MiOel.1}, the authors study the upper bound on $A[\tkernel]$. 
	Using the evolution equation satisfied in that work by the kernel $\kernel$ and hypotheses on its decay, 
	they prove that $A[\tkernel]$ is bounded uniformly in time. The bound on $A[\tkernel]$ allows then to fix the constant $C$. 
	So  the later term in the last inequality remains positive, permitting  $U$ to be a super-solution.
 	
 	As a simple example, if one computes $A[\tkernel]$ for $\kernel(a) = (1+a)^{4}$ one recovers:
 	$$
 	A[\kernel] = \frac{\int_0^\infty a \kernel(a+t)da }{\int_0^\infty \kernel(a+t)da} \sim \frac{t}{2}
 	$$
 	which forbids applying the same argument when the kernel $\kernel$ is constant in time with a tail of polynomial order.

 	Another argument should be used in the case of a compactly supported kernel. Indeed in this case there exists a time $t_0$ great enough s.t.
 	$\int_0^{t} \tkernel(a) da = 1$, which prevents from using \cite[Lemma 9.8.2]{grip} 
 \end{remark}
 
 
 \section{Numerical counterpart}\label{sec:discrete}
 In order to apply the previous ideas to the discrete setting, we first need to
 introduce a discrete framework that approximates the solutions of \eqref{eq.z}.
%  Namelly we discretize the time interval $[0,T]$ with a step $h>0$ and
%  we denote $t_n:= n h$ for all $n \in \{0, \ldots, N\}$ where $N:= T/h$. 
%  We assume for sake of simplicity that $T/h \in \N$.
 We use piecewise constant approximations. Indeed, for any function $g$ defined on $\rr$, we set:
 $ g^n := \frac{1}{h} \int_{t_n}^{t_{n+1}} g(t) dt$,  $\forall n \in \Z$.
 Using the rectangular rule, the discrete counterpart of \eqref{eq.z} reads: find $(Z^n)_{n\in\{0, \ldots, N\}}$ solving 
 \begin{equation}\label{eq.disc.z}
 	\left\{\begin{aligned}
 		&	Z^n - h \sum_{j\geq 0} Z^{n-j} \tr_j = f^n , & n \in \{0, \ldots, N\} \\
 		&	Z^n = Z_p^n , & n<0
 	\end{aligned}\right.
 \end{equation}
 where the discrete kernel $(\tr_j)_{j\in\N}$ is defined as in Section \ref{subsec:discrete} below.		
The sequence $(\tr_j)_{j\in\N}$ discretizes the kernel $\tkernel$:
 $
 R_j := \frac{1}{h} \int_{jh}^{(j+1)h} \kernel(a) da %, \quad \mu_{0,h} := h \sum_{j\geq 0} R_j, \quad \tr_j :=  R_j/ \mu_{0,h}.
 $.

\begin{lemma}
	Under the assumptions of Theorem \ref{thm.comparison}, the discrete kernel $(\tr_j)_{j\in\N}$ is a non-increasing sequence satisfying 
	$ - \frac{\tr_j - \tr_{j-1}}{h \tr_j} \leq \zeta$ for all $j \in \N^*$.
\end{lemma}

\begin{proof}
	Writing the difference : $R_{j+1}-R_j = \frac{1}{h} \int_{jh}^{(j+1)h} (\kernel(a+h) - \kernel(a)) da$, 
	then using the monotonicity of $\kernel$ provides that $R_{j+1} \leq R_j$ for all $j\in\N$. 
	Next, using again the monotonicity of $\kernel$ :
	\begin{align}
				R_{j+1}- R_{j} &= \frac{1}{h} \int_{jh}^{(j+1)h} (\kernel(a+h) - \kernel(a)) da 
				= \frac{1}{h} \int_{jh}^{(j+1)h} \int_a^{a+h} \da \kernel(s) ds da \\
				&\geq - \zeta \frac{1}{h} \int_{jh}^{(j+1)h} \int_a^{a+h} \kernel(s) ds da
				\geq - \zeta \frac{1}{h} \int_{jh}^{(j+1)h} h  \kernel(a+h)ds da 
				= -\zeta h R_{j+1}
	\end{align}		
	where we have used the that the kernel is non-increasing in the last inequality.
	Rearranging the latter inequality provides the desired result.			
\end{proof}


 \subsection{The initial layer}\label{subsec:discrete}
We solve:
 \begin{equation}\label{eq.disc.w}
 	\left\{\begin{aligned}
 		&	 W^n - h\sum_{j\geq 0} W^{n-j} \tr_j =0 & n \geq 0 \\
 		& W^n = - nh & n< 0.
 	\end{aligned}
 	\right.
 \end{equation} 
 The discrete elongation variable can be defined as in the continuous setting: $U^n_j := W^n - W^{n-j}$, and the analog of \eqref{eq.elongation.u} becomes
 \begin{equation}\label{eq.disc.elongation}
 	U^{n+1}_{j+1} = U^n_j + \delta W^{n+\nud} ,\quad \delta W^{n+\nud} := W^{n+1}-W^n, \quad j\geq0,\quad n\geq 0.
 \end{equation}
Expressing \eqref{eq.disc.w} with the elongation variable leads to the conservation principle $h \sum_{j\ge0} U^n_j \tr_j = 0$ for all $n\ge0$. 
		
\begin{lemma}
	The discrete finite differences satisfy:
	$$
	\delta W^{n+\nud} = - h \sum_{j\geq 1} {(\tr_j - \tr_{j-1})} U^{n+1}_j,\quad n\geq 0
	$$
\end{lemma} 
 
\begin{proof}
	Summing \eqref{eq.disc.elongation} against $h \tr_j$ for $j\geq 0$  
	%and adding $h U^{n+1}_1 \tr_0$ on both  sides
	provides: 
	$h \sum_{j\geq 0} \tr_j U^{n+1}_{j+1} = h \sum_{j\geq 0} \tr_j U^n_j + \delta W^{n+\nud}$, 
	then adding and subtracting $h \sum_{j\geq 0} \tr_j U^{n+1}_j$, one obtains:
	\begin{equation}\label{eq.elo.dw}
		h \sum_{j\geq 0} \tr_j U^{n+1}_j + h \sum_{j\geq 0} \tr_j (U^{n+1}_{j+1}- U^{n+1}_j) = h \sum_{j\geq 0} \tr_j U^n_j + \delta W^{n+\nud}.
	\end{equation}
	Performing an integration by parts transforms the second term on the left hand side into:
	$$
	h \sum_{j\geq 0} \tr_j (U^{n+1}_{j+1}- U^{n+1}_j)  = - h \sum_{j\geq 1} {(\tr_j - \tr_{j-1})} U^{n+1}_j
	$$
	and the conservation principle cancels the first terms on both sides of \eqref{eq.elo.dw}.
\end{proof} 
 
\begin{corollary}
	The elongation variable satisfies a closed equation:
	\begin{equation}\label{eq.closed.disc.elongation}
		U^{n+1}_{k+1} = U^n_k - h \sum_{j\geq 1} {(\tr_j - \tr_{j-1})} U^{n+1}_j, \quad k\geq 0 ,\quad n\geq 0
	\end{equation}
	which is the discrete version of \eqref{eq.elongation.u}.
\end{corollary} 
 
\begin{proposition}\label{prop.W.bounded}
	Assuming that there exists $\zeta>0$ s.t. for all $j \in \N^*$,  
	$$
	- \frac{R_j-R_{j-1}}{h R_j}  \leq \zeta 
	$$
	The discrete initial layer $(W^n)_{n\in\N}$ solving \eqref{eq.disc.w} is a bounded non-negative sequence (and the bound is uniform with respect to $h$).
\end{proposition} 
 
\begin{proof}[Proof of Proposition \ref{prop.W.bounded}]
 	We set $P^n_j:= h \sum_{\ell=0}^n U^\ell_j$, and we notice again that, as $\sum_{j\geq 0}\tr_j U^{\ell}_j \equiv 0$, 
	$\sum_{j\geq 0} P^n_j \tr_j \equiv 0$. Next using \eqref{eq.closed.disc.elongation}, one obtains that:
 	$$
 	P^{n+1}_{k+1} - h U^0_{k+1} = P^n_k - h\sum_{j\geq 1} (\tr_j-\tr_{j-1})  (P^{n+1}_j-h U^0_j), \quad k\geq 0 
 	$$
 	which is the discrete equivalent of \eqref{eq.p}. 
	Now we want to avoid constant source terms in the left hand side of 
	the previous expression so that we define $\tP_j^n := P^n_j - h \sum_{k\in\{0,\dots,j\}} U^0_k$  
	which transforms the previous equality into:
 	$$
 	\tP_{k+1}^{n+1} = \tP_{k}^n - h \sum_{j\geq 1} (\tr_j-\tr_{j-1}) 
	 \left(\tP^{n+1}_j+ h \sum_{\ell = 0}^{j-1}U^0_\ell\right)
 	$$
 	But $\sum_{j\geq 1}( \tr_j-\tr_{j-1}) \sum_{\ell=0}^{j-1} U^0_{\ell} 
	= - \sum_{j\geq 1} \tr_j U^0_j - \tr_0 U^0_0 = 0$, leading to 
 	$$
 	\tP_{k+1}^{n+1} = \tP_{k}^n - h \sum_{j\geq 1} (\tr_j-\tr_{j-1})  \tP^{n+1}_j %%- 2 h^2 \tr_0 U^0_0
 	$$

 	Applying the absolute value to the previous expression and integrating against $h \tr_k$ yields:
 	$$
 	I^n:= h \sum_{k\geq 0} \tr_k |\tP^{n+1}_{k+1}|  \leq h \sum_{k\geq 0} \tr_k |\tP_k^n|   
	- h \sum_{j\geq 1} (\tr_j-\tr_{j-1}) | \tP_j^n |  =: J^n
 	$$ 
 	But 
 	$I^n = h \sum_{k\geq 0} \tr_k |\tP^{n+1}_{k}| - h \sum_{k\geq 1}( \tr_k-\tr_{k-1}) |\tP^{n+1}_{k}| $, 
	and as $I^n \leq J^n$, this simplifies into:
 	$$
 	h \sum_{k\geq 0} \tr_k |\tP^{n+1}_{k}| \leq h \sum_{k\geq 0} \tr_k |\tP^{n}_{k}| %+ 2 h^2 \tr_0 \left|U^0_0\right| 
 	$$
 	and by induction this provides:
 	$$
 	h \sum_{k\geq 0} \tr_k |\tP^{n+1}_{k}|  
 	\leq h \sum_{k\geq 0} \tr_k |\tP^{0}_{k}| %+ 2 T h \tr_0 \left|U^0_0\right|  
 	\leq  h \sum_{k\geq 0} \tr_k \left|  h \sum_{\ell=0}^k U^0_\ell \right| %+ 2 h \frac{\mu_{1,h}}{\mu_{0,h}} 
 	\leq C(\mu_{2,h}, \mu_{1,h},\mu_0) 
 	$$
 	Returning to the definition of $(W^n)_{n\in\N}$,  one has 
 	$$
 	\delta W^{n+\nud} = - h \sum_{j\geq 1} (\tr_j-\tr_{j-1})U^{n+1}_j
 	$$
 	leading to 
 	$$
 	W^{n+1}-W^0 = -  \sum_{j\geq 1} (\tr_j-\tr_{j-1})(P^{n+1}_j- h U^0_j)
 	$$
 	and then 
 	$$
 \begin{aligned}
 		| W^{n+1}| \leq&  | W^0 | - h \sum_{j\geq 1} 
 		\frac{\tr_j-\tr_{j-1}}{h \tr_j} \tr_j \left( | P^{n+1}_j | +h | U_j^0 |\right) \\
 	& \leq | W^0 | + \zeta \left\{ h \sum_{j\geq 0}  \tr_j  | P^{n+1}_j | +h^2 \sum_{j\geq 0}  \tr_j    | U_j^0 |\right\}\\
 	& \leq | W^0 | + h \zeta \left\{ \sum_{j\geq 0}  \tr_j  | \tP^{n+1}_j | +h \sum_{j\geq 0}  \tr_j    \sum_{k=0}^{j} \left| U_k^0 \right| + c_1 \right\} \\
 	& \leq | W^0 | + h \zeta \left\{ 2 h \sum_{j\geq 0}  \tr_j    \sum_{k=0}^{j} \left| U_k^0 \right| + c_1 \right\} 	
 	\leq C(\mu_{0,h},\mu_{1,h},\mu_{2,h}) %\frac{\mu_{1,h}}{\mu_{0,h}} + 2 \zeta C \frac{\mu_{2,h}}{\mu_{1,h}}
 \end{aligned}
 	$$
Next we prove that $W^n \geq 0$ for all $n\geq 0$ by induction. 
First $W^0=\mu_{1,h}/(1-h \tr_0) \geq 0$ for $h$ 
small enough (and the bound is uniform by Lebesgue's Theorem for $h<h_0$). 
By definition of $W$, 
$W^{-n} =  -nh \geq 0$ for $n<0$. 
Assume now that $W^k \geq 0$ for all $k \leq n$, then using \eqref{eq.disc.w}, 
$
(1-h\tr_0) W^{n+1} = h \sum_{j\geq 1} \tr_j W^{n+1-j} \geq 0
$, 
which concludes the proof.
 \end{proof}
 
 \subsection{Discrete resolvent and comparison principle}

First we discretise \eqref{eq.z}. To do so we define a meshsize $h>0$ and define for sake of
simplicity the discrete kernel $(R_j)_{j\in\N}$ as in the previous section.


 By using successive convolutions of $(\tr_j)_{j\in\N}$, one can construct a non-negative $(Q_j)_{j\in\N}$ resolvent that solves 
 $$
 Q_n - h \sum_{j\geq 0}^n Q_{n-j} \tr_j = \tr_n, \quad \forall n \geq 1,\quad Q_0 = 0.
 $$
and we consider the problem \eqref{eq.disc.z}. We underline that we omit here 
the discretisation of the kernel on the first index $j=0$ in order to remain coherent 
with the quadrature formula of problem \eqref{eq.disc.z}. As the explicit value of the sequence $Q$
is not needed, this simplification does not alter the following results.

 The discrete convolution of two sequences $(U_j)_{j\in\N}$ and $(V_j)_{j\in\N}$ is defined as:
 $$
 (U \star V)_n := h \sum_{j=0}^n U_{n-j} V_j , \quad \forall n \geq 0.
 $$
 The solution of \eqref{eq.disc.z} can be expressed as $Z = f + f \star Q$.

\begin{theorem}\label{thm.comp.disc}
	Assume that $-(R_j-R_{j-1})/(hR_j) \leq \zeta$ for all $h>0$ and all $j\geq 1$, 
	and assumptions \ref{hypo.data}, then the solution of \eqref{eq.disc.z} 
	$(Z^n)_{n\in \{0,\dots,N\}}$ is bounded uniformly with respect to $h$.
\end{theorem}

 \begin{proof}[Proof of Theorem \ref{thm.comp.disc.main}]
 	First notice that thanks to the first hypothesis of the claim, 
	$h\sum_{j=0}^n \tr_j = 1 - h\sum_{j=n+1}^\infty  \tr_j \leq 1 - (1+\zeta h)^{-n+1} \tr_0/\zeta < 1$.
 	Then let's define $Q^m := \sum_{k=0}^m R^{k\star}$,  in the sense of sequences in $\ell^1$ such that 
	$(R\star R)_n = h\sum_{j=0}^n R_{n-j} R_j$ . 
 	One has that $Q = \lim_{m\to \infty} Q^m$ remains non-negative and finite for any $n\geq 0$. 
 	Indeed 
 	$|Q^m_n| \leq \sum_{k= 0}^m (h \sum_{j=0}^n R_j )^k \leq  \sum_{k= 0}^m  (1-\delta(n))^k \leq \delta(n)^{-1} < \infty$. The solution $Z$ can then be expressed as $Z=f+f \star Q$ 
 	and as the verification is purely algebraic one can follow the same steps as in \cite[Theorem 2.3.5 p. 44]{grip}. 
 	In the same spirit, one can extend \cite[Lemma 9.8.3]{grip} and show that if $S \geq S\star \tr + f$ 
	in the sense of convolution of sequences, then $S^n\geq Z^n$  for every $n\geq 0$, 
	where $Z$ is the solution of the equation $Z=\tr \star Z +f$.  

	The purpose is now to construct $S$. We set 
 	$$
 	S^n_f := \frac{|f|_\infty}{\mu_{1,h}} W^n + V^n , \quad V^n:=  C + \frac{1}{\mu_{1,h}}  \begin{cases}
 		h \sum_{j=0}^n \max_{k \in \{ j,N\}} |f^k| & \text{ if } n \geq 0 \\
 		h n |f|_\infty & \text{ otherwise}
 	\end{cases}
 	$$
 	then applying the integral operator gives: 
 	$$
 	\begin{aligned}
 		I^n &:=  S^n_f - h \sum_{j=0}^{n} S^{n-j}_f \tr_j \\
 		& \geq h \sum_{j=1}^n (V^n-V^{n-j})\tr_j 
 				+ h \sum_{j=n+1}^\infty (V^n-V^{n-j})\tr_j  
 				+  h \sum_{j\geq n+1}^\infty V^{n-j} \tr_j \\
 				& \qquad
 				+ h \frac{|f|_\infty}{\mu_{1,h}} \sum_{j=n+1}^\infty h(j-n) \tr_j\\
 		& \geq h \sum_{j=1}^n h \sum_{k=n-j}^n \max_{\ell \in \{ k,\dots,N\}} |f^\ell| \tr_j 
			+ h \sum_{j=n+1}^\infty \sum_{k=0}^n h \max_{\ell \in \{ k,\dots,N\}} |f^\ell| \tr_j
			+ \sum_{j=n+1}^\infty (V^0-V^{n-j} )\tr_j  \\
 		& \qquad+  h \sum_{j\geq n+1}^\infty V^{n-j} \tr_j  + h \frac{|f|_\infty}{\mu_{1,h}} \sum_{j=n+1}^\infty h(j-n) \tr_j\\
 		& \geq h \sum_{j=1}^n h j |f^n| \tr_j + h \sum_{j=n+1}^\infty \sum_{k=0}^n h n |f^n| \tr_j+ \frac{|f|_\infty}{\mu_{1,h}}\sum_{j=n+1}^\infty h({j-n} )\tr_j  \\
 		& \qquad+  h \sum_{j\geq n+1}^\infty V^{n-j} \tr_j  + h \frac{|f|_\infty}{\mu_{1,h}} \sum_{j=n+1}^\infty h(j-n) \tr_j \geq |f^n|\\
 	\end{aligned}
 	$$
 	The first line in the last inequality is then greater than $|f^n|$ and the second line cancels since the initial layer corrector's right hand side compensates exactly the
 	 tail of $U$ that was added and subtracted in the first lines above. One then compares:
 	 $$
 	 \left|Z^n\right|- h \sum_{j=0}^n \left| Z^{n-j} \right| \tr_j 
	 \leq 
	 |\tf^n|  \leq S_{\tf} - h \sum_{j=0}^n S^{n-j}_{\tf} \tr_j
 	 $$
 	where $\tf^n := f^n + h \sum_{j=n+1}^\infty |Z^{n-j}| \tr_j$, which implies that $|Z^n| \leq S^n$ as soon as $C \geq |Z^0|$.
 \end{proof}

\section{Application  to delayed sweeping processes}\label{sec:application}

We now apply the comparison principle developed in the previous sections to study delayed sweeping processes.
Recall Hypothesis \ref{hypo.convex.main} on the family of closed convex sets.

In this section we solve the following discrete problem: for any $n\geq 0$,
 \begin{equation}\label{eq.sweeping.projection}
	 X^n := P_{C^n}(\overline{X}^n), \quad \overline{X}^n := \frac{h \sum_{j\geq 1} X^{n-j} \tr_j}{h \sum_{j\geq 1} \tr_j}.
\end{equation}
where $P_{C^n}$ is the projection onto the convex set $C^n$ defined above. 

This problem can be reformulated as the following minimization problem: for any $n\geq 0$,
\begin{equation}\label{eq.sweeping}
	 X^n := \argmin_{W \in C^n} \cE_n(W), \quad \cE_n(W) := \frac{h}{2}\sum_{j\geq1} (W-X^{n-j})^2 \tr_j.
\end{equation}
 
\begin{proposition}[Minimization principle]
	The two minimizing problems \eqref{eq.sweeping.projection} and \eqref{eq.sweeping} are equivalent.
\end{proposition} 
 
\begin{proof}
	By definition of the projection onto a convex set: given $\overline{X}^n \in \mathbb{R}^d$,
	there exists a unique $X^n \in C^n$ s.t.
	$$
	X^n = \argmin_{W \in C^n} | W - \overline{X}^n |^2 
	$$
	Now we start from the energy functional $\cE_n$ and we compute:
	$$
	\begin{aligned}
		\cE_n(W) & = \frac{h}{2} \sum_{j\geq 1} 
		\left( W - \overline{X}^n + \overline{X}^n - X^{n-j} \right)^2 \tr_j \\
		& = \frac{h}{2} \sum_{j\geq 1} \left( W - \overline{X}^n \right)^2 \tr_j 
		+ h \sum_{j\geq 1} (W - \overline{X}^n)(\overline{X}^n - X^{n-j}) \tr_j	 \\
		& \quad + \frac{h}{2} \sum_{j\geq 1} (\overline{X}^n - X^{n-j})^2 \tr_j 
		= (1-h \tr_0)| W - \overline{X}^n |^2 +  R((X^{n-j})_{j\geq 1})
	\end{aligned}
	$$
	where the remainder term $R$ does not depend on $W$. This concludes the proof.
\end{proof}

 
%  The Euler Lagrange equations (in their Kuhn and Tucker version)  associated to the minimization problem \eqref{eq.sweeping} read :
%  \begin{equation}\label{eq.sweep.eul}
%  	X^n - \overline{X}^n + \lambda^n \varphi_n' (X^n)=0, \quad \overline{X}^n:= h \sum_{j\geq 1} X^{n-j} \tr_j
%  \end{equation}
%  or reformulated in the context of projectors on convex sets:
%  \begin{equation}\label{eq.}
%  	X^n = P_{C^n}(\overline{X}^n), \quad n\geq 0
%  \end{equation}
 
% \begin{proposition}\label{prop.4.1}
% 	Under assumptions above, and assuming that the past condition $X_p^{-1}\in C^0$, $X$ the unique discrete solution of the minimization problem \eqref{eq.sweeping} admits a uniform BV bound with respect to the discretization parameter $h$.
% \end{proposition}

% \begin{proof}
% We start with the following inequality:
% $$
% \begin{aligned}
% 	| X^n - &X^{n-1} |  = \left|P_{C^n}(\overline{X}^n) - P_{C^{n-1}}(\overline{X}^{n-1})\right| \\
% 	& \leq \left|P_{C^n}(\overline{X}^n) - P_{C^{n-1}}(\overline{X}^{n})\right|  
% 		 + \left|P_{C^{n-1}}(\overline{X}^n) - P_{C^{n-1}}(\overline{X}^{n-1})\right| \\
% 	&  \leq d_H(C^{n-1},C^n)  + \left| \overline{X}^n - \overline{X}^{n-1}\right|  
% 	   \leq d_H(C^{n-1},C^n)  + h\sum_{j\geq 1} \left| X^{n-j}- X^{n-j-1}\right| \tr_j \\
% 	&  \leq  d_H(C^{n-1},C^n)  + h\sum_{j= 1}^n  	   \left|X^{n -j}- X^{n-j-1}\right| \tr_j 
% 							  + h\sum_{j= n+1}^\infty  \left|X^{n -j}- X^{n-j-1}\right| \tr_j
% \end{aligned}
% $$ 
% where in the first line we have used the triangle inequality, in the second 
% line the non-expansiveness of the projection operator and in the third line the definition of $\overline{X}^n$.		
% As $x_p \in \Lip(\mathbb{R}^-)$ with the Lipschitz constant $L_{x_p}$, the past condition contributes as:
% $$
% h\sum_{j= n+1}^\infty  \left|  {X}^{n -j}- {X}^{n-1-j}\right| \tr_j\leq L_{x_p} h
% $$
% then we denote by $y^n:= \sum_{k=1}^n 	| X^k - X^{k-1} |$ and write:
% $$
% \begin{aligned}
% 	y^n  &  \leq \TV(C) + \sum_{k=1}^n   h\sum_{j= 1}^k  \left| {X}^{k -j}- {X}^{k-1-j}\right| \tr_j +  n h L_{x_p}  \\
% 	&  \leq \TV(C) + h\sum_{j= 1}^n  y^{n-j}\tr_j+  h \sum_{j=1}^n |X^0-X^{-1}| \tr_j+  n h L_{x_p}  \\
% 	&  \leq \TV(C) + h\sum_{j= 1}^n  y^{n-j}\tr_j+  \mu_{0,h} |X^0-X^{-1}| + T L_{x_p} 
% \end{aligned}
% $$
% hereafter we detail the transition from the first to the second line:
% $$
% \begin{aligned}
% 	& h \sum_{k=1}^n  \sum_{j= 1}^k  \left| {X}^{k -j}- {X}^{k-1-j}\right| \tr_j  
% 	= h\sum_{j=1}^n  \sum_{k= j}^n  \left| {X}^{k -j}- {X}^{k-1-j}\right| \tr_j  
% 	= h \sum_{j=1}^n  \sum_{k= 0}^{n-j}  \left| {X}^{k}- {X}^{k-1}\right| \tr_j \\
% 	&= h\sum_{j=1}^n  \sum_{k= 1}^{n-j}  \left| {X}^{k}- {X}^{k-1}\right| \tr_j 
% 		+ h\sum_{j=1}^n    \left| {X}^{0}- {X}^{-1}_p \right| \tr_j= h  \sum_{j=1}^n y^{n-j} \tr_j 
% 		+ h\sum_{j=1}^n    \left| {X}^{0}- {X}^{-1}_p \right| \tr_j 
% \end{aligned}
% $$
% The first step of the discrete  differences contains a jump:
% $$
% \begin{aligned}
% 	|X^0-X^{-1}_p | \leq &  \left| P_{C^0}(\overline{X}^0) -P_{C^0}(X_p^{-1})\right| \leq |\overline{X}^0-X^{-1}_p | \leq h \sum_{j=0}^\infty |X^{-j}_p -X^{-1}_p| \tr_j \\
% 	& \leq L_{x_p} h \sum_{j\ge0}   hj \tr_j \leq L_{x_p}   \frac{\mu_{1,h}}{\mu_{0,h}}
% \end{aligned}
% $$
% We can summarize the previous estimates into:
% $$
% y^n \leq  h\sum_{j= 1}^n  y^{n-j}\tr_j+  \gamma 
% $$
% where $\gamma:= \TV(C) + C(\mu_0,\mu_1,L_{x_p}) T$. Then one remarks that for $h$ sufficiently small,
% $$
% (1-h \tr_j) y^n \leq  y_n  \leq  h\sum_{j= 1}^n  y^{n-j}\tr_j+  \gamma \implies y^n- \sum_{j=0}^{n} y^{n-j}\tr_j \leq \gamma
% $$
% the latter estimate being compatible with the discrete comparison principle framework.
% Then applying Theorem \ref{thm.comp.disc.main}, one concludes that $y^n \leq T \{ \TV(C) + C(\mu_0,\mu_1,L_{x_p})\}$ which ends the claim.
% \end{proof}

% \begin{remark}
% 	 It does not seem  possible to provide  a finer result following ideas from \cite{Moreau}. Indeed, 
% 	 assuming that $X_p^k \in C^0$ for $k<0$, for any $n\geq 0$, 
% 	$$
% 	\begin{aligned}
% 		|X^{n+1}-X^n| \leq&  |X^{n}-\overline{X}^n| + |\overline{X}^n-\overline{X}^{n-1}| + |\overline{X}^{n-1}-X^{n-1}| \\
% 		& = \inf_{y \in C^{n}}  |y-\overline{X}^n| +\inf_{y \in C^{n-1}}  |y-\overline{X}^{n-1}| +
% 		|\overline{X}^n-\overline{X}^{n-1}|   \\
% 		& 	\leq h \sum_{j\geq 1}\inf_{y \in C^{n}}  |y-X^{n-j}|  \tr_j +  h\sum_{j\geq 1} \inf_{y \in C^{n-1}}  |y-X^{n-1-j}|  \tr_j 
% 		+ |\overline{X}^n-\overline{X}^{n-1}|   \\
% 		& \leq  h \sum_{j\geq 1}^n e(C^{n-j},C^n)  \tr_j + h \sum_{j\geq 1}^{n-1} e(C^{n-1-j},C^{n-1})  \tr_j  \\
% 		& \qquad + h \sum_{j\geq n+1} e(C^{0},C^n)  \tr_j + h \sum_{j\geq n} e(C^{0},C^{n-1})  \tr_j  + |\overline{X}^n-\overline{X}^{n-1}|   \\
% 	\end{aligned}
% 	$$
% 	Next if we assume  that $C(t)$ is of finite retraction, {\em i.e.}: by considering
% 	all the finite sequences $(\tau_i)_i$ of the form $s = \tau_0  < \tau_1< \dots \leq \tau_N = t$
% 	$$
% 	\sup_{(\tau_i)_i} \sum_{i=1}^{N_{(\tau_i)_i}} e(C(\tau_{i+1}),C(\tau_i)) =: \ret(C; s,t) = r(t)-r(s)
% 	$$
% 	where $e(A,B):= \sup_{a\in A} \inf_{b\in B} |a-b|$, the function $r$ is non-decreasing finite on any segment $[s,t]\subset [0,T]$, which leads to:
% 	$$
% 	\begin{aligned}
% 		|X^{n+1}-X^n| \leq&    h \sum_{j\geq 1}^n (r(nh)-r(n-j)h)  \tr_j + h \sum_{j\geq 1}^{n-1} (r((n-1)h)-r((n-1-j)h)) \tr_j  \\
% 		& \qquad +r(nh)  h \sum_{j\geq n+1} \tr_j + r((n-1)h) h \sum_{j\geq n}  \tr_j  + |\overline{X}^n-\overline{X}^{n-1}| \\
% 		& \leq 4 |r|_\infty + |\overline{X}^n-\overline{X}^{n-1}| 
% 	\end{aligned}
% 	$$
% 	which is not summable. This comes from the fact that we express retraction with respect to all past positions of the set $C^n$ 
% 	which increases too drastically the upper  bound.
% \end{remark}


% \begin{lemma}\label{lem.proj}
% 	Let A and B be two closed convex sets in $\mathbb{R}^d$, then for any $x\in \mathbb{R}^d$,
% 	$$
% 	| P_A(x) - P_B(x) | \leq d_H(A,B)
% 	$$
% \end{lemma}


% \begin{proof}
% We proceed by a double inequality.
% For any $a\in A$, one has there exists $b\in B$ s.t. $|a-b| \leq d_H(A,B)$, then
% $$
% d(x,B) \leq d(x,b) \leq |x-a| + |a-b| \leq |x-a| + d_H(A,B)
% $$
% Taking the infimum over $a\in A$ leads to
% $$
% d(x,B) \leq d(x,A) + d_H(A,B)
% $$
% Interchanging the roles of $A$ and $B$ gives the reverse inequality and then
% $$
% | d(x,A) - d(x,B) | \leq d_H(A,B)
% $$
% Now by definition of the projection onto a convex set, one has:
% $$
% | P_A(x) - x | = d(x,A), \quad | P_B(x) - x | = d(x,B)
% $$
% which concludes the proof.
% \end{proof}

 
 \subsection{Energy estimates}
 
\begin{proposition}\label{prop.nrj}
	Assume that the sequence $(\tr_j)_{j\in \N}$ is decreasing and that the convex sets satisfy
	Hypothesis \ref{hypo.convex.main}. 
	Then the discrete sweeping sequence $(X^n)_{n\ge0}$ defined by \eqref{eq.sweeping}	 
	satisfies the energy estimate:	
	\[
	\cE_{N}(X^{N}) +  \sum_{n=0}^{N-1} D_n \le e^{2 T} \cE_{0}(X^0) + e^{2 T} \sum_{n=0}^{N-1} e_n,
	\]
	where $D_n= \frac{h}{2} \sum_{j\ge1} (\tr_{j}-\tr_{j+1}) | X^{n} - X^{n-j}|^2 \geq 0$ is the dissipation term and
	$e_n:= e(C^n,C^{n+1})^2/h$.
\end{proposition}
 
 \begin{proof}[Proof of Proposition \ref{prop.nrj.main}]
 	The minimization principle \eqref{eq.sweeping} gives
 	\[
 	\cE_{n+1}(X^{n+1}) \le \cE_{n+1}(P_{C^{n+1}}(X^n)).
 	\]
 	Adding and subtracting $X^n$ inside the energy functional leads to
 	\[
 	\begin{aligned}
 		\cE_{n+1}(X^{n+1})
 		& \le \frac{1}{2}  
 		| X^n - P_{C^{n+1}}(X^n)|^2
 		+ \sum_{j\ge1}( X^n - P_{C^{n+1}}(X^n),\,X^{n+1-j}-X^n)\tr_j \\
 		& + \frac{h}{2} \sum_{j\geq 1} | X^{n+1-j} - X^n|^2 \tr_j \\
 	\end{aligned}
 	\]
 	We denote $e_n:= e(C^n,C^{n+1})^2/h$, and write accordingly
 	\[
 	\begin{aligned}	
 		\cE_{n+1}(X^{n+1})
 		& \le \frac{h}{2} e_n 
 		+ h \sum_{j\ge1} e(C^n,C^{n+1})|X^{n+1-j}-X^n|\tr_j
 		+  \frac{h}{2} \sum_{j\geq 1} | X^{n+1-j} - X^n|^2 \tr_j .
 	\end{aligned}
 	\] 
 	Using Young's inequality on the second term on the right-hand side gives
 	\[
 	\begin{aligned}	
 		\cE_{n+1}(X^{n+1})
 		& \le (1+2 h ) e_n 
 		+ (1+2 h) \cE_{n+1}(X^n).		
 	\end{aligned}
 	\]
 	Then using the definition of the energy functional and rearranging terms leads to
 	\[	
 	\cE_{n+1}(X^n) = \cE_{n}(X^n) - \frac{h}{2} \sum_{j\ge1} (\tr_{j+1}-\tr_j) | X^{n} - X^{n-j}|^2 =: \cE_{n}(X^n) - D_n,
 	\]
 	where $D_n= \frac{h}{2} \sum_{j\ge1} (\tr_{j}-\tr_{j+1}) | X^{n} - X^{n-j}|^2 \geq 0$ is the dissipation term.
 	Putting everything together, we obtain
 	\[
 	\cE_{n+1}(X^{n+1}) + (1+2 h) D_n \le (1+2 h) \cE_{n}(X^n) + (1+2 h ) e_n.
 	\]
 	Iterating this inequality from $n=0$ to $n=N-1$ gives
 	\[
 	\cE_{N}(X^{N}) + (1+2 h) \sum_{n=0}^{N-1} D_n \le (1+2 h)^N \cE_{0}(X^0) + (1+2 h ) \sum_{n=0}^{N-1} (1+2 h)^{N-1-n} e_n.
 	\]
 	Since $(1+2 h)^N \le e^{2 T}$ and $\sum_{n=0}^{N-1} (1+2 h)^{N-1-n} e_n \le e^{2 T} \sum_{n=0}^{N-1} e_n$, the proof is complete.	
 \end{proof}
 
 \begin{remark}
 	The energy $\cE_{N}(X^N)$ completes the dissipation term at step $n=N$ since one writes:
 	$$
 	h \sum_{p=0}^N \sum_{j\geq1} (\tr_j - \tr_{j+1}) \left| X^{p}- X^{p-j}\right|^2 \leq 2 \cE_{N}(X^{N}) + 2 h\sum_{p=0}^{N-1} \sum_{j\geq1} (\tr_j - \tr_{j+1}) \left| X^{p}- X^{p-j}\right|^2  
 	$$
 \end{remark}
 
\begin{corollary}\label{coro.dissip}
	Under the previous hypotheses, 
	$$
	\sum_{p= 0}^N h 
	\left( \sum_{j\ge1} (\tr_{j-1} - \tr_{j}) \left| X^{p} - X^{p-j}\right| \right)^2 \leq  C < \infty 
	$$
	uniformly with respect to the discretization step $h$.
\end{corollary} 
 
\begin{proof}%[Proof of Corollary \ref{coro.dissip.main}]
	Since the sequence $(\tr_{j})_{j\geq 0}$ is non-increasing and of mass 1 (indeed $h \sum_{j\geq 0} \tr_{j}=1$), one has that 
	$$
	\sum_{j\ge1} \tr_{j-1} - \tr_{j} = \tr_0 \leq (h \Lip_\kernel/2  + \kernel(0))/\mu_{0} < \infty 
	$$
 where we have used the Lipschitz hypothesis in Assumptions \ref{hypo.data}. Then using Cauchy-Schwartz, 
	one writes:
	$$
	\begin{aligned}
		\sum_{p= 0}^N & h \left( \sum_{j\geq1} (\tr_{j-1} -\tr_{j}) |X^p-X^{p-j}| \right)^2  \leq \sum_{p=0}^N h \sum_{j\geq 1} (\tr_{j-1} -\tr_{j})
								\sum_{j\geq 1} (\tr_{j-1} -\tr_{j}) |X^p-X^{p-j}|^2\\
		& \leq  \tr_0 \sum_{p= 0}^N h \sum_{j\geq 1}  (\tr_{j-1} -\tr_{j}) |X^p-X^{p-j}|^2 
		\leq  \tr_0 \sum_{p= 0}^N D^p < \infty
	\end{aligned}
	$$
\end{proof} 


\begin{corollary}\label{coro.lag}
	Under the previous hypotheses, for $h$ small enough, one has the uniform bound:
	$$
	\left|\overline{X}^n - X^n\right| \leq C 
	$$
\end{corollary}

\begin{proof}%[Proof of Corollary \ref{coro.lag}]
	Using Cauchy-Schwartz inequality, one writes:
	$$
	\begin{aligned}
		\left|\overline{X}^n - X^n\right|^2 & 
		= \left| \frac{h \sum_{j\geq 1} (X^{n-j} - X^n) \tr_j}{1-h \tr_0} \right|^2  
		\leq  \frac{h \sum_{j\geq 1} (X^{n-j}-X^n)^2 \tr_j }{1-h \tr_0} \times \frac{h \sum_{j\geq 1} \tr_j}{1-h \tr_0} 		\leq \frac{2 \cE_n(X^n)}{(1-h \tr_0)} < \infty \\
	\end{aligned}
	$$
	where we have used the energy estimate provided by Proposition \ref{prop.nrj.main}.	
\end{proof}
\begin{remark}
	We notice that contrarily to the standard sweeping process without delay, 
	the distance between the projected point $X^n$ and the unconstrained point $\overline{X}^n$ 
	remains only bounded whereas in the standard case it is $h$ close. 
	This greatly complicates the 
	analysis and does not allow to obtain straightforward $BV$ estimates in time as in \cite{mor77}. 
\end{remark}

% \subsection{Case of convex sets}

% \textit{Under construction!}

% Cutting straight to the point, we aim to prove a small enough lower bounds for
% $$
% \langle \lambda^{n+1} G^{n+1} - \lambda^n G^n, X^{n+1} - X^n \rangle
% $$
% where $G^n \in \partial \varphi_n(X^n)$ and $\lambda^n \geq 0$ are the Lagrange multipliers associated to the minimization problem \eqref{eq.sweeping}.
 
% We notice that $X^{n+1} \in C^{n+1}$ and $X^n \in C^n$, but it can happen that $X^{n+1} \notin C^n$ and $X^n \notin C^{n+1}$. Otherwise, the monotonicity of normal cones of convex sets would directly imply that this term is non-negative.

% \subsubsection{Estimates using the Hausdorff-distance}

% Our first aim is to find nearby points $Y^{n+1} \in C^n$ and $Y^n \in C^{n+1}$ such that we can use convexity to write:
% $$
% \langle \lambda^{n+1} G^{n+1}, X^{n+1} - Y^n \rangle \geq 0, \quad	
% \langle \lambda^n G^n, X^{n} - Y^{n+1} \rangle \geq 0.
% $$

% For this, we define the projection:
% $$
% Y^{n+1}:= P_{C^n}(X^{n+1})
% $$

% Due to the assumptions \eqref{eq.b.c}, we have the estimate:
% $$
% |X^{n+1} - Y^{n+1}| = d_{C^{n}}(X^{n+1}) \leq d_H(C^{n+1},C^n) \leq C h.
% $$

% Now we can compute:
% $$
% \langle G^{n}, X^{n+1} - X^n \rangle = \langle G^{n}, X^{n+1} - Y^{n+1} \rangle + \langle G^{n}, Y^{n+1} - X^n \rangle \leq \langle G^{n}, X^{n+1} - Y^{n+1} \rangle
% $$

% Using Cauchy-Schwartz, we get:
% $$
% \langle G^{n}, X^{n+1} - X^n \rangle \geq - |G^n| |X^{n+1} - Y^{n+1}| \geq - C h |G^n|
% $$

% Similar, we obtain with switched roles of $n$ and $n+1$:
% $$\langle G^{n+1}, X^{n+1} - X^n \rangle \leq C h |G^{n+1}|.$$

% Combined, we have:
% $$
% \langle \lambda^{n+1} G^{n+1} - \lambda^n G^n, X^{n+1} - X^n \rangle \leq C h ( \lambda^{n+1} |G^{n+1}| + \lambda^n |G^n|).
% $$

% % remark 
% \begin{remark}
% 	We notice that $G^n$ and $X^{n+1} - Y^{n+1}$ are typically expected to be pointing in similar directions (or be zero), but without further assumptions it can happend that this contribution is negative, for example if $C^{n} = [0,\frac{h}{4}] + n h$ and $X^0 = \frac{h}{4}$ and $X^1 = h$, then $\langle G^1 - G^0, X^1 - X^0 \rangle = - \frac{h}{4}$. This shows that the provided estimate is sharp in general.
% \end{remark}

% \subsubsection{Estimates using regularity of constraint functions}

% We can obtain sharper estimates if we assume regularity of the constraint function $\varphi$.

\subsection{The case of a moving circular set}\label{sec:circular}
 
In this section we prove Theorem \ref{thm:circular.main}. We consider the case where the convex sets are moving circles:
 \begin{equation}
 	C(t):= \overline{B}(c(t), R) = \{x \in \mathbb{R}^m:  \phi_t(x):= |x - c(t)|^2 -	 R^2 \leq 0\}
 \end{equation}
 with $c \in H^1([0,T];\mathbb{R}^m)$ and $R \in \mathbb{R}_+$.
We discretize the sets as $C^n:= C(nh)$ for $n=0,\dots,N$ with $N = \lfloor T/h \rfloor$.
The corresponding discrete sweeping process is defined by \eqref{eq.sweeping}. Under these assumptions,
Hypothesis \ref{hypo.convex.main} is satisfied since
 $$ d_H(C^{n},C^{n-1}) \leq |c(nh)-c((n-1)h)| $$ 
 and $c$ being in $H^1([0,T];\mathbb{R}^m)$ implies that 
 $$
 \sum_{n=1}^{N} d_H(C^{n},C^{n-1})^2/h \leq \sum_{n=1}^{N} |c(nh)-c((n-1)h)|^2/h \leq |c|_{H^1([0,T];\mathbb{R}^m)}^2 < \infty.
 $$
We denote by $\varphi_n(X):= |X - c^n|^2 - R^2$ where $c^n:= c(nh)$. These constraints
are qualified in the sense of \cite{Ciarlet89} since $\nabla \varphi_n(X) = 2 (X - c^n) \neq 0$ 
for all $X$ such that $\varphi_n(X) =0$.	This provides the necessary condition
in order to  write Euler Lagrange equations (in their Kuhn and Tucker version)  
associated to the minimization problem \eqref{eq.sweeping}, they read:
 \begin{equation}\label{eq.sweep.eul}
 	h \sum_{j\geq 0} (X^n - X^{n-j})\tr_j  + \lambda^n \nabla \varphi_n (X^n)=0, \quad \forall n\geq 0.
 \end{equation}
 

\begin{proof}[Proof of Theorem \ref{thm:circular.main}]
 	We recall the definition of $\phi^n(X):= |X - c^n|^2 - R^2$ where $c^n:= c(nh)$.
 	From the Euler-Lagrange equation \eqref{eq.sweep.eul}, setting the elongation variable as 
 	$U^n_j:= X^n-X^{n-j}$ for $j\geq 1$, it  satisfies by definition 
 	$$
 	U^{n+1}_{j+1} = U^n_j + X^{n+1} - X^n 
 	$$
 	which summed against $\tr_j$ gives
 	\begin{equation}\label{eq:inner.prod}
 		h \sum_{j\geq 0} U^{n+1}_{j+1} \tr_j = h \sum_{j\geq 0} U^n_j \tr_j + \mu_{0} (X^{n+1}-X^n) = 
 		h \sum	_{j\geq 0} U^n_j \tr_j + X^{n+1}-X^n , 
 	\end{equation}
 	since $\mu_0 = h \sum_j \tr_j =1$. 
 	Adding and subtracting $U^{n+1}_j$ in the left hand side gives
 	$$
 	h \sum_{j\geq 0} U^{n+1}_{j+1} \tr_j = h \sum_{j\geq 0} U^{n+1}_j \tr_j + h \sum_{j\geq 0} (U^{n+1}_{j+1}-U^{n+1}_j) \tr_j 
 	= h \sum_{j\geq 1} U^{n}_j \tr_j + (X^{n+1}-X^n).					
 	$$
 	By a change of indices one recovers that $  \sum_{j\geq 1} (U^{n+1}_{j+1}-U^{n+1}_j) \tr_j $ $= \sum_{j\geq 1} 
 	U^{n+1}_j(\tr_{j-1}-\tr_{j}) =: d^{n+1}$,  where $d^{n}$ is related to the dissipation term $D^n$
 	provided by Proposition \ref{prop.nrj.main} (to be precised below).
 	The Euler-Lagrange equation \eqref{eq.sweep.eul},
 	can be rephrased as 
 	$$
 	h \sum_{j\ge 0 } U^{n}_j \tr_{j}+ \lambda^n G^n = 0
 	$$
 	and provides using the equality above:
 	$$
 	-\lambda^{n+1} G^{n+1} +  d^{n+1} =  -\lambda^n G^n + (X^{n+1}-X^n)
 	$$	
 	where we defined $G^n:= \nabla \varphi_n (X^n) = 2(X^n - c^n)$.
 	We test the previous equality by $X^{n+1}-X^n$:
 	$$
 	|X^{n+1}-X^n|^2 + \left( \lambda^{n+1}  G^{n+1} - \lambda^n G^n, X^{n+1}-X^n \right) 
 	=  h \left( d^{n+1}, X^{n+1}-X^n \right)
 	$$
 	and we consider the second term on the left-hand side:
 	$$
 	\begin{aligned}
 		2 & \left(\lambda^{n+1}  G^{n+1} - \lambda^n G^n, X^{n+1}-X^n \right)  = 
 		 \left(\lambda^{n+1}  G^{n+1} - \lambda^n G^n, G^{n+1}-G^n \right) + 2 \left(\lambda^{n+1}  G^{n+1} - \lambda^n G^n, c^{n+1}-c^n \right) \\
 		& = \frac{1}{2} \left( \lambda^{n+1}+ \lambda^{n}\right) | G^{n+1}-G^n |^2 
 		+ \frac{1}{2} (\lambda^{n+1}- \lambda^{n}) ((G^{n+1})^2 - (G^n)^2) 
 		+
 		2 \left(\lambda^{n+1}  G^{n+1} - \lambda^n G^n, c^{n+1}-c^n \right)	 \\
 		& = \frac{1}{2} \left( \lambda^{n+1}+ \lambda^{n}\right) | G^{n+1}-G^n |^2 
 		+ 2 (\lambda^{n+1}- \lambda^{n}) (\phi^{n+1}(X^{n+1}) - \phi^n(X^n)) \\
 		& \quad +
 		2\left(\lambda^{n+1}  G^{n+1} - \lambda^n G^n, c^{n+1}-c^n \right)
 	\end{aligned}		
 	$$
 	Notice first that since $(\lambda^n)_{n\in \N}$ is non-negative, the first term is non-negative and we can neglect it in the following estimates. 
 	We focus now on the second term. 
 	\begin{itemize}[itemsep=0.05cm,parsep=0.1cm]
 		\item Assume that the constraint is active at step $n$ and $n+1$, i.e. $\lambda^n >0$ and $\lambda^{n+1} >0$. 
 		Then $\phi^n(X^n) = 0$ and $\phi^{n+1}(X^{n+1})= 0$ and the second term vanishes.
 		\item Assume that the constraint is inactive at step $n$ and $n+1$, i.e. $\lambda^n =0$ and $\lambda^{n+1} =0$. 
 		Then the second term vanishes as well.
 		\item Assume that the constraint is active at step $n$ and inactive at step $n+1$, i.e. 
 		$\lambda^n >0$, $\phi^n(X^n) = 0$	 and $\lambda^{n+1} =0$, $\phi^{n+1}(X^{n+1}) < 0$, 
 		the second term reduces to $-\lambda^n \phi^{n+1}(X^{n+1}) \geq 0$.
 		\item Assume that the constraint is inactive at step $n$ and active at step $n+1$, i.e. 
 		$\lambda^n =0$, $\phi^n(X^n) < 0$ and $\lambda^{n+1} >0$, $\phi^{n+1}(X^{n+1}) = 0$. 
 		the second term reduces to $\lambda^{n+1} (-\phi^n(X^n)) \geq 0$.
 	\end{itemize}
 	In all cases, the second term is non-negative.
 	To sum up, one has
 	$$
 	\begin{aligned}
 		(X^{n+1}-X^n)^2 & \leq - \langle \lambda^{n+1} G^{n+1} - \lambda^n G^n, c^{n+1}-c^n \rangle 
 		+  h \langle d^{n+1}, X^{n+1}-X^n \rangle
 	\end{aligned}	
 	$$
 	Using then that $\lambda^{n+1} G^{n+1} - \lambda^n G^n = - (X^{n+1}-X^n) + h \sum_{j\geq 0} (X^{n+1-j}-X^{n-j}) \tr_j$
 	from \eqref{eq.sweep.eul}, 
 	denoting $\delta X^{n+\nud} = X^{n+1}-X^n$, we have	using Young's inequality:
 	$$
 	\begin{aligned}			
 		|\delta X^{n+\nud}|^2 & \leq \left( \delta X^{n+\nud}, c^{n+1}-c^n \right) 
 		- \left( h \sum_{j\geq 0} \delta {X}^{n+\nud-j} \tr_j, c^{n+1}-c^n \right) 
		+ h \left( d^{n+1}, \delta X^{n+\nud} \right) \\
 		& \leq |\delta X^{n+\nud}| |c^{n+1}-c^n| 
		+ h \left|\sum_{j\geq 0} \delta {X}^{n+\nud-j} \tr_j\right| |c^{n+1}-c^n| 
 		+  h |d^{n+1}| |\delta X^{n+\nud}| \\
 		& \leq \alpha_1 |\delta X^{n+\nud}|^2 + \alpha_1^{-1} |\delta c^{n+\nud}|^2 
 		+ \alpha_2 \left|h \sum_{j\geq 0} \delta {X}^{n+\nud-j} \tr_j\right|^2  + \alpha_2^{-1} |c^{n+1}-c^n|^2 \\
 		& \qquad+ h^2  \alpha_3^{-1} |d^{n+1}|^2 +  \alpha_3 |\delta X^{n+\nud}|^2
 	\end{aligned}
 	$$
 	Leading to the final estimate:
 	$$
 	\begin{aligned}			
 		|\delta X^{n+\nud}|^2 &\leq 
 		(\alpha_1^{-1} + \alpha_2^{-1}) /(1 - \alpha_1 - \alpha_3)  |c^{n+1}-c^n|^2 \\
 		& + \alpha_2/(1 - \alpha_1 - \alpha_3)  \left|h \sum_{j\geq 0} \delta {X}^{n+\nud-j} \tr_j\right|^2 + h	 \alpha_3 D^{n+1} 	/(1 - \alpha_1 - \alpha_3	) 		
 	\end{aligned}
 	$$
 	One chooses $\alpha_2=1-\alpha_1 - \alpha_3$ and dividing the previous expression by $h$ gives:
 	$$
 	\begin{aligned}			
 		y^n:= & \frac{|\delta X^{n+\nud}|^2}{h} \leq 
 		\frac{(\alpha_1^{-1} + (1-\alpha_1 - \alpha_3)^{-1})}{(1 - \alpha_1 - \alpha_3)}  \frac{|c^{n+1}-c^n|^2}{h}  
		+ \frac{1}{h} \left|h \sum_{j\geq 0} \delta {X}^{n+\nud-j} \tr_j\right|^2 + 	 \frac{\alpha_3}{(1 - \alpha_1 - \alpha_3	)} |D^{n+1}|^2	\\
 		& \leq c \omega_n  + h \sum_{j\geq 0} y^{n-j} \tr_j 	
 	\end{aligned}
 	$$
 	where $\omega_n:= (c^{n+1}-c^n)^2/h + D^{n+1} $ and we used Jensen's inequality to 
 	estimate $\left|h \sum_{j\geq 0} \delta {X}^{n+\nud-j} \tr_j\right|^2$.
 	Define now $z^n:= \sum_{k=0}^{n} y^k$, then summing the previous inequality gives:
 	$$
 	z^n \leq c \sum_{k=0}^n \omega_k + h \sum_{j\geq 0}^{n} \tr_j z^{n-j}
 	$$
 	We detail here how to obtain the last sum: for $n \geq 1$ 
	%{\color{red}very delicate to be checked because of the jump at $t=0$}.
 	$$
 	%\begin{aligned}
 		\sum_{p= 0}^{n}\sum_{j\geq 0} y^{p-j} \tr_j 
		= \sum_{p=0}^{n} \sum_{j=0}^p y^{p-j} \tr_j = \sum_{0\leq j \leq p \leq n} y^{p-j} \tr_j 
		= \sum_{j= 0}^{n} \sum_{k=0}^{n-j} y^{k} \tr_j 
 		= \sum_{j=0}^{n} z^{n-j} \tr_j
 	%\end{aligned}
 	$$
 	since we assumed that the data is well prepared, for all $j \geq 1$ $X^{-j} = X^{-1}\in C^0$ implying that $y^k = 0$ for $k < 0$.

	
 	Using then Theorem \ref{thm.comp.disc.main}  leads to the desired estimate since 
 	$\sum_{k=0}^n \omega_k \leq C$. Indeed, 
	since $\dt c \in L^2 (0,T)$, its equivalent 
 	on the discrete level provides $\sum_{n=0}^{N_T} (\delta c^{n+ \nud})^2/h < \infty$. Moreover we use also 
 	the estimate provided by Corollary \ref{coro.dissip}. This ends the proof.	
 \end{proof}
 
\begin{remark}
	The key step of the  proof relies on the sign of 
	$\langle\lambda^{n+1}  G^{n+1} - \lambda^n G^n, X^{n+1}-X^n \rangle$ which has a 
	positive part on one side and a component that depends on the trajectories of the 
	set's centers on the right hand side.
	This heavily depends on the geometry of the moving set. To our knowledge, it is an open
	problem in the general case. One does not see how to extend these results as 
	a function, for instance, of the quantity $e(C^{n},C^{n+1})$ used in the previous 
	section to establish energy and dissipation estimates.   
\end{remark} 
 
 
 \subsection{Convergence towards the delayed sweeping process}
 
 
 \begin{proof}[Proof of Theorem \ref{thm:limit_sweeping.main}]
 	\emph{Step~1 -- Compactness of $\{X_h\}_h$.}
 	From the uniform $H^1$ bound, we obtain (Kondrachov’s selection theorem) that there exists
 	$X\in H^1([0,T];\mathbb R^m)$ and a subsequence such that
 	$X_h\to X$ in $C^0(0,T)$.
	Moreover, the piecewise-constant interpolant $\tilde{X}_h$ also converges to $X$ in $L^p(0,T)$.
 	
 	\emph{Step~2 -- Convergence of the discrete averages.}
	By Proposition \ref{prop:discrete_convolution_convergence}, 
	the discrete convolution defining $\overline{X}_h$ converges to the continuous convolution $\overline{X}$
	in $L^p$ for any $1\le p<\infty$. The continuous limit convolution decomposes accordingly as
 	\[
 	\overline{X}(t)
 	%=\mathcal A_{\mathrm{int}}(t)+\mathcal A_{\mathrm{past}}(t)
 	:=\int_0^t X(t-s)\varrho(s)\,ds
 	+ X_p \int_t^\infty \varrho(s)\,ds.
 	\]
 	
 	\emph{Step~3 -- Convergence in the discrete inclusion.}
 	At the discrete level, for each $n\ge0$, we have the projection characterization $
 	X^n = P_{C^n}(\overline{X}^n) $ is equivalent to the normal cone inclusion
 	$\overline{X}^n - X^n \in N_{C^n}(X^n)$ or, in duality form,
 	\[
	\forall \xi \in \R^m, \langle \overline{X}^n - X^n , \xi \rangle \leq \left| \overline{X}^n - X^n \right| d_{C^n}(\xi + \overline{X}^n - X^n).
 	\]
 	Define the piecewise-constant sets $C_h(t):=C^n$ for $t\in(nh,(n+1)h]$.
 	By the $H^1(0,T)$ assumption on the centers of balls  $c(t)$,
 	we may assume $C_h(t)\to C(t)$ in the Hausdorff sense for a.e.\ $t$.
 	Moreover by Corollary~\ref{coro.lag}, we have that $X^n-\overline{X}^n$ is uniformly bounded in $\R^m$.

	Up to a subsequence, we may fix $t\in(0,T)$ such that
 	$\overline{X}_h(t)\to \overline{X}(t)$, $\tilde{X}_h(t)\to X(t)$,
 	and $C_h(t)\to C(t)$.
 	Then for such $t$,
 	\[
 	\langle \overline{X}_h(t) - \tilde{X}_h(t) , \xi \rangle 
 	\leq \left| \overline{X}_h(t) - \tilde{X}_h(t) \right| d_{C_h(t)}(\xi + \overline{X}_h(t) - \tilde{X}_h(t)).
 	\]
	Because one has that 
	$$
	\begin{aligned}
		\left| d_{C_h(t)} (\xi + \right. & \overline{X}_h(t)- \tilde{X}_h(t))  \left.- d_{C(t)}(\xi + \overline{X}(t) - X(t))\right| \\
	& \leq \left| (\overline{X}_h(t) - \tilde{X}_h(t))- (\overline{X}(t) - X(t)) \right| + d_H(C_h(t),C(t)),
	\end{aligned}
	$$
	and that $d_{C_h(t)} (\xi + \overline{X}_h(t) - \tilde{X}_h(t))$ is uniformly bounded for every fixed $\xi$ 
	thanks to the boundedness of $\overline{X}_h(t) - \tilde{X}_h(t)$ (see Corollary~\ref{coro.lag}),
 	we obtain, by successive triangle inequalities and passing to the limit $h\to0$,
 	\[
 	\langle \overline{X}(t) - X(t) , \xi \rangle 
 	\leq \left| \overline{X}(t) - X(t) \right| d_{C(t)}(\xi + \overline{X}(t) - X(t)).
 	\]
 	Since this holds for all $\xi\in\R^m$, which is equivalent to the normal cone inclusion
 	$
 	\overline{X}(t)-X(t)\in N_{C(t)}\big(X(t)\big)
 	$.	
 \end{proof}
 
 
%
 

\appendix


	 \section{Discrete approximation of kernels' moments}\label{sec:moments}
	
	In order to compare the continuous quantities 
	\(\mu_k:= \int_0^{\infty} a^k \varrho(a)\,da\) 
	with their discrete counterparts that appear in the estimates above,
	we introduce for every \(h>0\) the discrete moments
	\[
	\mu_{k,h}:= h \sum_{j\ge0} (jh)^k R_j,
	\qquad 
	R_j:= \frac{1}{h}\int_{jh}^{(j+1)h}\varrho(a)\,da,
	\qquad k\in\{0,1,2\}.
	\]
	Note that by construction the discrete mass is exact:
	\[
	\mu_{0,h}=h\sum_{j\ge0}R_j=\int_0^\infty\varrho(a)\,da=\mu_0.
	\]
	The following lemma shows that these quantities are first–order accurate
	approximations of the continuous moments under our integrability assumptions.
	
	\begin{lemma}[Approximation of the moments]
		\label{lem:moments}
		Assume that the kernel satisfies \(\varrho\ge0\) and 
		\(\int_0^{\infty}(1+a)^{k-1}\varrho(a)\,da<\infty\) 
		for \(k\in\{1,2\}\). Then there exist constants \(C_1,C_2>0\),
		independent of \(h\in(0,1]\), such that
		\begin{equation}\label{eq:moment_error}
			|\mu_k-\mu_{k,h}|
			\;\le\;
			C_k\,h
			\int_0^{\infty}(1+a)^{k-1}\varrho(a)\,da,
			\qquad
			C_1=1,\; C_2=3.
		\end{equation}
	\end{lemma}
	
	\begin{proof}
		For every \(k\ge1\),
		\[
		\mu_k
		=\sum_{j\ge0}\int_{jh}^{(j+1)h}a^k\varrho(a)\,da.
		\]
		For \(a\in[jh,(j+1)h]\), by the mean value theorem,
		\(
		|a^k-(jh)^k|
		\le k (jh+h)^{k-1} h.
		\)
		Hence
		\[
		|\mu_k-\mu_{k,h}|
		\le
		k h \sum_{j\ge0}(jh+h)^{k-1}
		\int_{jh}^{(j+1)h}\varrho(a)\,da
		\le
		C_k h\int_0^{\infty}(1+a)^{k-1}\varrho(a)\,da,
		\]
		where \(C_k=k2^{k-1}\), giving the claimed constants
		\(C_1=1\) and \(C_2=3\).
	\end{proof}
	
	
	\begin{remark}
		Estimates \eqref{eq:moment_error}
		justify the substitution of discrete moments
		\(\mu_{i,h}\) for their continuous counterparts
		\(\mu_i\) in all bounds of discrete {\em a priori} estimates above. 
		In particular, the constants appearing in those inequalities
		remain uniform in \(h\) as \(h\to0\).
	\end{remark}
	

	
\section{$L^p$ Convergence of Discrete Convolutions}\label{sec:discrete_convolution}

Let $T > 0$ and $h = T/N$ for $N \in \mathbb{N}$. Define the intervals $J_n = [nh, (n+1)h)$ for $n=0, \dots, N-1$. 
Consider the piecewise constant functions:
\[ B_h(t) = \sum_{n=0}^{N-1} B_n \chi_{J_n}(t), \quad C_h(t) = \sum_{n=0}^{N-1} C_n \chi_{J_n}(t) \]

\begin{proposition}\label{prop:discrete_convolution_convergence}
Assume $B_h \to B$ strongly in $L^p(0,T)$ and $C_h \to C$ strongly in $L^1(0,T)$. 
Define the discrete causal convolution $A_h$ as:
\[ A_h(t) = \sum_{n=0}^{N-1} \left( h \sum_{j=0}^n B_{n-j} C_j \right) \chi_{J_n}(t) \]
Then $A_h \to B * C$ strongly in $L^p(0,T)$, where $(B * C)(t) = \int_0^t B(t-s)C(s)ds$.
\end{proposition}

\begin{proof}
The proof proceeds in two main steps: identifying $A_h$ as a projection of a continuous convolution, and applying Young's Inequality.

\paragraph{Step 1: Identification of the Operator}
For piecewise constant functions $B_h$ and $C_h$, the continuous convolution $(B_h * C_h)(t)$ evaluated at a grid point $t = nh$ is:
\[ (B_h * C_h)(nh) = \int_0^{nh} B_h(nh-s)C_h(s)ds = \sum_{j=0}^{n-1} \int_{jh}^{(j+1)h} B_h(nh-s)C_h(s)ds \]
Since $B_h$ and $C_h$ are constant on each interval $J_j$, the integral becomes $h B_{n-1-j} C_j$. Thus, the discrete sum in $A_h$ corresponds exactly to the values of the continuous convolution of the piecewise constant approximations at the nodes. Specifically, $A_h$ is the piecewise constant interpolant of $(B_h * C_h)$.

\paragraph{Step 2: Decomposition of the Error}
By the triangle inequality in $L^p(0,T)$:
\[ \|A_h - B * C\|_{L^p} \leq \|A_h - B_h * C_h\|_{L^p} + \|B_h * C_h - B * C\|_{L^p} \]

The first term $\|A_h - B_h * C_h\|_{L^p}$ vanishes as $h \to 0$ because $B_h * C_h$ is a continuous (piecewise linear) function, and $A_h$ is its piecewise constant approximation. 

For the second term, we apply Young's Inequality for convolutions ($\|f * g\|_{L^p} \leq \|f\|_{L^p}\|g\|_{L^1}$):
\[ \|B_h * C_h - B * C\|_{L^p} \leq \|(B_h - B) * C_h\|_{L^p} + \|B * (C_h - C)\|_{L^p} \]
\[ \leq \|B_h - B\|_{L^p}\|C_h\|_{L^1} + \|B\|_{L^p}\|C_h - C\|_{L^1} \]

Since $B_h \to B$ in $L^p$, the term $\|B_h - B\|_{L^p} \to 0$. Since $C_h \to C$ in $L^1$, the term $\|C_h - C\|_{L^1} \to 0$. By the stability of strong convergence, $\|C_h\|_{L^1}$ is uniformly bounded. Therefore, both terms approach zero.
And the claim follows.
\end{proof}

\section{Decay of a decreasing functions in weighted \(L^1\) spaces}\label{sec:decay_lipschitz}
\begin{proposition}\label{prop:decay_weighted}
Let $\kernel: \mathbb{R}_+ \to \mathbb{R}$ be a monotone non increasing non-negative function such that 
$\kernel \in L^1(\mathbb{R}_+, (1+a)^2)$. 
Then $\lim_{a \to \infty} \kernel(a)(1+a)^2 = 0$.
\end{proposition}

\begin{proof}
	Let's $f(a) = \kernel(a)(1+a)^2$. By hypothesis, $f \in L^1(\mathbb{R}_+)$.
	Denote $f^n:= \frac{1}{h} \int_{nh}^{(n+1)h} f(a) da$ for $n \in \mathbb{N}$ and $h > 0$.
	Since $f$ is non-negative and integrable, we have $\lim_{n \to \infty} f^n = 0$.
	Moreover, since $\kernel$ is non-increasing,
	\[
	f^{n-1}\geq \kernel(nh) \frac{1}{h}  \int_{(n-1)h}^{nh} (1+a)^2 da \geq  
	\kernel(nh) (nh)^2.
	\]
	for any $n \geq 1$. Thus one has:
	$$
	(1+a)^2 \kernel(a) \leq 2 f^{n-1} \quad \text{ for } a \in [nh, (n+1)h]
	$$
	which shows the desired limit by letting $n \to \infty$.
\end{proof}



	
%\end{appendices}
% \section{Internal-part approximation (contribution $1\le j\le n$)}
% 
% Recall the notation: for $t\in[0,T]$ set $n=\lfloor t/h\rfloor$ and
% \[
% \mathcal A^h_{\mathrm{int}}(t)
%:=h\sum_{j=1}^{n} R_j\,\overline X^h(t-jh),
% \qquad
% \mathcal A_{\mathrm{int}}(t)
%:=\int_0^{t} X(t-s)\,\widetilde\varrho(s)\,ds,
% \]
% with $hR_j=\int_{jh}^{(j+1)h}\varrho(s)\,ds$ and $\widetilde\varrho=\varrho/\mu_0$.
% 
% \begin{lemma}[Convergence of the internal part]\label{lem:internal_part}
% 	Assume $\varrho\ge0$, $\mu_0=\int_0^\infty\varrho>0$, and the cell averages $R_j$ are
% 	mass-preserving. Let $\{\overline X^h\}_h$ be the right-continuous
% 	piecewise-constant interpolants of the discrete solutions and assume
% 	\(\overline X^h\to X\) strongly in \(L^1(0,T;\mathbb R^m)\) as \(h\to0\).
% 	Assume also $X\in L^1_{\mathrm{loc}}(\mathbb R_+;\mathbb R^m)\).
% 	Then
% 	\[
% 	\mathcal A^h_{\mathrm{int}}\longrightarrow \mathcal A_{\mathrm{int}}
% 	\quad\text{in }L^1(0,T;\mathbb R^m).
% 	\]
% \end{lemma}
% 
% \begin{proof}
% 	Fix $T'>T$ (or take any $T'<\infty$ larger than $T$) and let $S>0$ be a truncation parameter.
% 	Write for $t\in[0,T]$ (with $n=\lfloor t/h\rfloor$)
% 	\[
% 	\mathcal A^h_{\mathrm{int}}(t)-\mathcal A_{\mathrm{int}}(t)
% 	= I_h(t) + J_h(t),
% 	\]
% 	where
% 	\[
% 	I_h(t):=h\sum_{j=1}^{n} R_j\big(\overline X^h(t-jh)-X(t-jh)\big),
% 	\]
% 	\[
% 	J_h(t):=h\sum_{j=1}^{n} R_j X(t-jh)-\int_0^t X(t-s)\widetilde\varrho(s)\,ds.
% 	\]
% 	
% 	\emph{Estimate for $I_h$.} Integrate in $t$ on $[0,T]$ and change the summation/integration order:
% 	\[
% 	\int_0^T |I_h(t)|\,dt
% 	\le \int_0^T h\sum_{j\ge1} R_j \big|\overline X^h(t-jh)-X(t-jh)\big|\,dt
% 	= \sum_{j\ge1} hR_j \int_0^T \big|\overline X^h(t-jh)-X(t-jh)\big|\,dt.
% 	\]
% 	Perform the change of variable $\tau=t-jh$ in each inner integral (note that for $j\le n$ the integration range remains included in $(-S,T')$ for a fixed finite $S$ if $h$ small enough), which yields
% 	\[
% 	\int_0^T |I_h(t)|\,dt
% 	\le \mu_0 \,\|\overline X^h - X\|_{L^1(-S,T')}.
% 	\]
% 	Since $\overline X^h\to X$ strongly in $L^1(0,T)$ and is uniformly bounded on compact intervals (by BV), we can let $h\to0$ and then let $S\to0$ (or simply choose $S$ fixed large enough and use strong $L^1$-convergence on a slightly bigger compact interval) to conclude $\|I_h\|_{L^1(0,T)}\to0$.
% 	
% 	\emph{Estimate for $J_h$.} Split the integral and discrete sum at a finite truncation $S>0$:
% 	\[
% 	\begin{aligned}
% 		J_h(t) = &\underbrace{\Big(h\sum_{1\le j\le \lfloor S/h\rfloor} R_j X(t-jh)
% 			- \int_0^S X(t-s)\widetilde\varrho(s)\,ds\Big)}_{=:J_h^{(1)}(t)}\\
%& 		+
% 		\underbrace{\Big(h\sum_{j>\lfloor S/h\rfloor}^{n} R_j X(t-jh)
% 			-\int_S^{t} X(t-s)\widetilde\varrho(s)\,ds\Big)}_{=:J_h^{(2)}(t)}.
% 	\end{aligned}
% 	\]
% 	For fixed $S$ the first term $J_h^{(1)}(t)$ is a classical quadrature error for the convolution of the fixed function $X$ against the integrable kernel $\widetilde\varrho$ on $[0,S]$. Using the cell-average definition of $R_j$ (mass-preserving quadrature) and dominated convergence, we get $J_h^{(1)}\to0$ in $L^1(0,T)$ as $h\to0$.
% 	
% 	For the second term $J_h^{(2)}(t)$ we bound its $L^1(0,T)$-norm by the tail contribution:
% 	\[
% 	\int_0^T |J_h^{(2)}(t)|\,dt
% 	\le
% 	\int_0^T \Big(h\sum_{j>\lfloor S/h\rfloor}^{\infty} R_j |X(t-jh)|\Big) dt
% 	+
% 	\int_0^T \int_S^\infty |X(t-s)|\widetilde\varrho(s)\,ds\,dt.
% 	\]
% 	Using Fubini/Tonelli and the integrability of $\widetilde\varrho$, the second integral equals
% 	\[
% 	\int_S^\infty \widetilde\varrho(s) \int_0^T |X(t-s)|\,dt\,ds
% 	\le \int_S^\infty \widetilde\varrho(s) \int_{-S}^{T} |X(\tau)|\,d\tau\,ds,
% 	\]
% 	which can be made arbitrarily small by choosing $S$ large (uniformly in $h$) because $\widetilde\varrho\in L^1$ and $X\in L^1_{\mathrm{loc}}$. The discrete tail term is controlled similarly by replacing the integral over the cell by the cell mass:
% 	\[
% 	\int_0^T h\sum_{j>\lfloor S/h\rfloor} R_j |X(t-jh)|\,dt
% 	= \sum_{j>\lfloor S/h\rfloor} hR_j \int_0^T |X(t-jh)|\,dt
% 	\le \Big(\sum_{j>\lfloor S/h\rfloor} hR_j\Big)\int_{-S}^{T}|X(\tau)|\,d\tau,
% 	\]
% 	and the prefactor $\sum_{j>\lfloor S/h\rfloor} hR_j=\int_S^\infty\varrho(s)\,ds$ tends to $0$ as $S\to\infty$, uniformly in $h$.
% 	
% 	Putting together the estimates for $J_h^{(1)}$ and $J_h^{(2)}$, we first choose $S$ large to make the tails small and then let $h\to0$ to make the quadrature error small; hence $\|J_h\|_{L^1(0,T)}\to0$.
% 	
% 	Combining the two parts gives $\|\mathcal A^h_{\mathrm{int}}-\mathcal A_{\mathrm{int}}\|_{L^1(0,T)}\to0$.
% \end{proof}
% 
% \begin{remark}[Rates under extra regularity]
% 	If, in addition, $X$ is uniformly continuous on compact intervals with modulus $\omega_X(\cdot)$,
% 	one obtains the quantitative bound
% 	\[
% 	\|\mathcal A^h_{\mathrm{int}}-\mathcal A_{\mathrm{int}}\|_{L^1(0,T)}
% 	\le \mu_0\,\|\overline X^h-X\|_{L^1(-S,T')}\;+\;\int_0^T \big(\omega_X(h)+o(1)\big)\,dt,
% 	\]
% 	so that if $\overline X^h$ is constructed by cell-averaging $X$ (or is piecewise-linear converging
% 	with order $h$) one recovers an $O(h)$ convergence rate.
% \end{remark}
% 
% 
% 
% 
% 
% \section{Discrete approximation of the past tail}
% 
% We keep the notation of the main text. For \(t\in[0,T]\) and \(h>0\) write \(n=\lfloor t/h\rfloor\)
% and recall the discrete past contribution
% \[
% \mathcal A^h_{\mathrm{past}}(t)
%:=h\sum_{j\ge n+1} R_j\,X_p(t-jh),
% \qquad\text{with}\qquad hR_j=\int_{jh}^{(j+1)h}\varrho(s)\,ds,
% \]
% and the continuous past tail
% \[
% \mathcal A_{\mathrm{past}}(t)
%:=\int_t^\infty X_p(t-s)\,\widetilde\varrho(s)\,ds,
% \qquad \widetilde\varrho=\frac{\varrho}{\mu_0},\ \mu_0=\int_0^\infty\varrho.
% \]
% 
% \begin{lemma}[Discrete-to-continuous approximation of the past tail]\label{lem:past_tail}
% 	Assume \(\varrho\ge0\), \(\mu_0=\int_0^\infty\varrho>0\), and \(R_j\) are the mass-preserving
% 	cell averages defined above. Let \(X_p:(-\infty,0]\to\mathbb{R}^m\) be Lipschitz on \((-\infty,0]\)
% 	with constant \(L\). Then, for every \(t\in[0,T]\),
% 	\begin{equation}\label{eq:past_uniform}
% 		\big|\mathcal A^h_{\mathrm{past}}(t)-\mathcal A_{\mathrm{past}}(t)\big|
% 		\le L\,h.
% 	\end{equation}
% 	As a consequence,
% 	\begin{equation}\label{eq:past_L1}
% 		\|\mathcal A^h_{\mathrm{past}}-\mathcal A_{\mathrm{past}}\|_{L^1(0,T)}
% 		\le L\,h\,T.
% 	\end{equation}
% 	If instead \(X_p\) has modulus of continuity \(\omega_{X_p}(\cdot)\) on \((-\infty,0]\),
% 	then \eqref{eq:past_uniform} and \eqref{eq:past_L1} hold with \(Lh\) replaced by \(\omega_{X_p}(h)\).
% \end{lemma}
% 
% \begin{proof}
% 	Fix \(t\in[0,T]\) and put \(n=\lfloor t/h\rfloor\). Using \(hR_j=\int_{jh}^{(j+1)h}\varrho(s)\,ds\)
% 	and writing the continuous tail as a cellwise integral we obtain
% 	\[
% 	\mathcal A_{\mathrm{past}}(t)
% 	= \sum_{j\ge n+1}\int_{jh}^{(j+1)h} X_p(t-s)\,\frac{\varrho(s)}{\mu_0}\,ds,
% 	\qquad
% 	\mathcal A^h_{\mathrm{past}}(t)
% 	= \sum_{j\ge n+1} \Big(\int_{jh}^{(j+1)h}\frac{\varrho(s)}{\mu_0}\,ds\Big)\, X_p(t-jh).
% 	\]
% 	Hence for each cell \(j\ge n+1\) the local contribution to the pointwise error is
% 	\[
% 	\Big|\int_{jh}^{(j+1)h}\frac{\varrho(s)}{\mu_0}\big(X_p(t-jh)-X_p(t-s)\big)\,ds\Big|
% 	\le \int_{jh}^{(j+1)h}\frac{\varrho(s)}{\mu_0}\,|X_p(t-jh)-X_p(t-s)|\,ds.
% 	\]
% 	Because \(s\in[jh,(j+1)h]\) and \(t-jh<0\) for \(j\ge n+1\), the arguments \(t-jh\) and \(t-s\)
% 	both belong to \((-\infty,0]\). Using the Lipschitz property of \(X_p\),
% 	\(|X_p(t-jh)-X_p(t-s)|\le L\,|s-jh|\le Lh\). Thus the integrand is bounded by \(Lh\)
% 	and integrating against the probability weight \(\varrho(s)/\mu_0\) on the cell yields
% 	the per-cell bound \(Lh\int_{jh}^{(j+1)h}\varrho(s)/\mu_0\,ds\). Summing over \(j\ge n+1\)
% 	gives
% 	\[
% 	|\mathcal A^h_{\mathrm{past}}(t)-\mathcal A_{\mathrm{past}}(t)|
% 	\le Lh \sum_{j\ge n+1}\int_{jh}^{(j+1)h}\frac{\varrho(s)}{\mu_0}\,ds
% 	= Lh\int_{(n+1)h}^\infty\frac{\varrho(s)}{\mu_0}\,ds
% 	\le Lh,
% 	\]
% 	which proves \eqref{eq:past_uniform}.
% 	
% 	Inequality \eqref{eq:past_L1} follows immediately by integrating the uniform in \(t\)
% 	bound over \(t\in[0,T]\).
% 	
% 	If \(X_p\) is not Lipschitz but has modulus of continuity \(\omega_{X_p}(\cdot)\),
% 	replace the bound \(L|s-jh|\) above by \(\omega_{X_p}(|s-jh|)\le\omega_{X_p}(h)\),
% 	and the same steps give the variant with \(\omega_{X_p}(h)\) in place of \(Lh\).
% \end{proof}
% 
% \begin{remark}
% 	The estimate is sharp in the sense of order: the sampling error on each cell is proportional
% 	to the cell size \(h\) times the local variation of \(X_p\). If one uses a higher-order
% 	reconstruction of the past values (e.g. mid-point sampling together with extra regularity
% 	of \(X_p\)), the error can improve to \(O(h^2)\) or better.
% \end{remark}
%% 
%
% 
\bibliographystyle{abbrv}
\bibliography{biblio}


\end{document}
Isabelle Dupin CRCTB, données biologiques sur la secretion du mucus
Mathieu Delorme CRCTB, impact de l'air sur dynamique du mucus



s

